\documentclass[12pt,a4paper,oneside]{book}

% ============================================================================
% PACKAGES
% ============================================================================
\usepackage[utf8]{inputenc}
\usepackage[T1]{fontenc}
\usepackage{mathptmx} % Times New Roman font
\usepackage{geometry}
\usepackage{setspace}
\usepackage{graphicx}
\usepackage{listings}
\usepackage{xcolor}
\usepackage{amsmath}
\usepackage{amssymb}
\usepackage{hyperref}
\usepackage{fancyhdr}
\usepackage{titlesec}
\usepackage{tocloft}
\usepackage{tabularx}
\usepackage{booktabs}
\usepackage{multirow}
\usepackage{float}
\usepackage{enumitem}
\usepackage{caption}
\usepackage{subcaption}
\usepackage{pmboxdraw} % Support for Unicode box-drawing characters

% ============================================================================
% PAGE GEOMETRY
% ============================================================================
\geometry{
    a4paper,
    left=1.5in,
    right=1in,
    top=1in,
    bottom=1in
}

% ============================================================================
% LINE SPACING
% ============================================================================
\onehalfspacing

% ============================================================================
% HEADERS AND FOOTERS (IEEE Format)
% ============================================================================
\pagestyle{fancy}
\fancyhf{}
\fancyhead[L]{\leftmark}
\fancyhead[R]{\thepage}
\renewcommand{\headrulewidth}{0.4pt}
\setlength{\headheight}{15pt}

% Remove date from plain page style (affects second page)
\fancypagestyle{plain}{%
    \fancyhf{}%
    \fancyfoot[C]{\thepage}%
    \renewcommand{\headrulewidth}{0pt}%
}

% ============================================================================
% IEEE-STYLE TITLE FORMATTING (Bold)
% ============================================================================
% Part formatting - IEEE style with bold
\titleformat{\part}[display]
{\normalfont\huge\bfseries\centering}
{\partname\ \thepart}
{20pt}
{\Huge\bfseries}

% Chapter formatting - IEEE style with bold
\titleformat{\chapter}[display]
{\normalfont\LARGE\bfseries}
{\chaptertitlename\ \thechapter}
{20pt}
{\Huge\bfseries}

% Section formatting - IEEE style with bold
\titleformat{\section}
{\normalfont\Large\bfseries}
{\thesection}
{1em}
{\Large\bfseries}

% Subsection formatting - IEEE style with bold
\titleformat{\subsection}
{\normalfont\large\bfseries}
{\thesubsection}
{1em}
{\large\bfseries}

% Subsubsection formatting - IEEE style with bold
\titleformat{\subsubsection}
{\normalfont\normalsize\bfseries}
{\thesubsubsection}
{1em}
{\normalsize\bfseries}

% ============================================================================
% CAPTION FORMATTING (IEEE Style - Bold)
% ============================================================================
\captionsetup{
    labelfont=bf,
    textfont=bf,
    justification=centering,
    singlelinecheck=false,
    font=small
}

% ============================================================================
% TABLE OF CONTENTS FORMATTING (Bold)
% ============================================================================
\renewcommand{\cftpartfont}{\bfseries\Large}
\renewcommand{\cftchapfont}{\bfseries}
\renewcommand{\cftsecfont}{\bfseries}

% ============================================================================
% HYPERREF SETUP
% ============================================================================
\hypersetup{
    colorlinks=true,
    linkcolor=blue,
    filecolor=magenta,
    urlcolor=cyan,
    citecolor=green,
    pdftitle={A1M: A 64-bit RISC Architecture with Complete Compiler Toolchain},
    pdfauthor={Shahrear Hossain, Riazul Islam},
}

% ============================================================================
% CODE LISTING SETUP - Professional Black and White Theme (Enhanced/Bold)
% ============================================================================
\definecolor{codebg}{rgb}{0.95,0.95,0.95}        % Slightly darker gray background
\definecolor{codeframe}{rgb}{0.2,0.2,0.2}        % Darker frame for contrast
\definecolor{codekeyword}{rgb}{0.0,0.0,0.0}      % Black keywords (bold)
\definecolor{codecomment}{rgb}{0.25,0.25,0.25}   % Darker gray comments
\definecolor{codestring}{rgb}{0.15,0.15,0.15}    % Darker strings
\definecolor{codenumber}{rgb}{0.3,0.3,0.3}       % Darker line numbers

% Base listing style - Professional black and white with enhanced boldness
\lstdefinestyle{base}{
    backgroundcolor=\color{codebg},
    commentstyle=\color{codecomment}\itshape,
    keywordstyle=\color{codekeyword}\bfseries,
    numberstyle=\tiny\color{codenumber},
    stringstyle=\color{codestring},
    basicstyle=\ttfamily\footnotesize\fontseries{b}\selectfont,
    breakatwhitespace=false,
    breaklines=true,
    captionpos=b,
    keepspaces=true,
    numbers=left,
    numbersep=8pt,
    showspaces=false,
    showstringspaces=false,
    showtabs=false,
    tabsize=4,
    frame=single,
    frameround=tttt,
    rulecolor=\color{codeframe},
    framerule=0.8pt,
    xleftmargin=2em,
    framexleftmargin=1.5em,
    aboveskip=1.5em,
    belowskip=1.5em
}

% Verilog style
\lstdefinestyle{verilog}{
    style=base,
    language=Verilog,
    morekeywords={module,endmodule,input,output,wire,reg,always,posedge,
                  negedge,parameter,localparam,begin,end,if,else,case,
                  endcase,default,assign,initial,for}
}

% Assembly style
\lstdefinestyle{assembly}{
    style=base,
    language=[x86masm]Assembler,
    morekeywords={MOVI,ADD,ADDI,SUB,SUBI,MUL,MULI,DIV,DIVI,MOD,MODI,
                  AND,ANDI,OR,ORI,XOR,XORI,SHL,SHLI,SHR,SHRI,SAR,SARI,
                  LOAD,STORE,LOADB,STOREB,JMP,JZ,JNZ,JGT,JLT,JGE,JLE,
                  CALL,RET,CMP,CMPI,PUSH,POP,SYSCALL,NOP,HALT,MOV},
    comment=[l]{;},
    morecomment=[l]{//}
}

% C style
\lstdefinestyle{c}{
    style=base,
    language=C,
    morekeywords={uint64_t,uint32_t,uint16_t,uint8_t,int64_t,int32_t,
                  int16_t,int8_t,bool,size_t,NULL}
}

% Algo style
\lstdefinestyle{algo}{
    style=base,
    morekeywords={fun,return,if,else,while,for,int,bool,char,long,
                  true,false,print,print_int,print_char,read_int,read_char},
    comment=[l]{//},
    morecomment=[s]{/*}{*/},
    morestring=[b]",
    morestring=[b]'
}

% Bash/Shell style
\lstdefinestyle{bash}{
    style=base,
    language=bash,
    morekeywords={echo,cd,ls,make,gcc,./,sudo,install}
}

% ============================================================================
% CHAPTER AND SECTION FORMATTING
% ============================================================================
\titleformat{\chapter}[display]
{\normalfont\huge\bfseries}{\chaptertitlename\ \thechapter}{20pt}{\Huge}
\titlespacing*{\chapter}{0pt}{0pt}{40pt}

\titleformat{\section}
{\normalfont\Large\bfseries}{\thesection}{1em}{}
\titlespacing*{\section}{0pt}{3.5ex plus 1ex minus .2ex}{2.3ex plus .2ex}

\titleformat{\subsection}
{\normalfont\large\bfseries}{\thesubsection}{1em}{}
\titlespacing*{\subsection}{0pt}{3.25ex plus 1ex minus .2ex}{1.5ex plus .2ex}

% ============================================================================
% CUSTOM COMMANDS
% ============================================================================
\newcommand{\code}[1]{\texttt{#1}}
\newcommand{\keyword}[1]{\textbf{\code{#1}}}
\newcommand{\register}[1]{\texttt{R#1}}
\newcommand{\instruction}[1]{\texttt{#1}}
\newcommand{\hex}[1]{\texttt{0x#1}}

% Professional box for important notes
\newcommand{\notebox}[2]{
\begin{center}
\fbox{\parbox{0.9\textwidth}{
\textbf{#1:} #2
}}
\end{center}
}

% ============================================================================
% BEGIN DOCUMENT
% ============================================================================
\begin{document}

% ============================================================================
% TITLE PAGE
% ============================================================================
\begin{titlepage}
    \centering
    \vspace*{1cm}
    
    {\LARGE\bfseries International Islamic University Chittagong\par}
    \vspace{0.5cm}
    {\large Department of Electrical and Electronics Engineering\par}
    \vspace{2cm}
    
    {\Huge\bfseries A1M: A 64-bit RISC Architecture\\[0.5cm]
    with Complete Compiler Toolchain\par}
    \vspace{1cm}
    
    {\Large\itshape A Comprehensive Study of Computer Architecture,\\
    Assembly Language, and Compiler Design\par}
    \vspace{2cm}
    
    {\Large\bfseries Submitted By:\par}
    \vspace{0.5cm}
    {\large Shahrear Hossain Shawon\par}
    \vspace{2cm}
    
    {\Large\bfseries Supervised By:\par}
    \vspace{0.5cm}
    {\large Lecturer Riazul Islam\par}
    {\large Department of Electrical and Electronics Engineering\par}
    {\large International Islamic University Chittagong\par}
    \vspace{2cm}
    
    \vfill
    
    {\small This work is submitted in partial fulfillment of the requirements\\
    for research in Computer Architecture and Compiler Design\par}
\end{titlepage}

% ============================================================================
% DECLARATION
% ============================================================================
\chapter*{Declaration}
\addcontentsline{toc}{chapter}{Declaration}

I hereby declare that the work presented in this thesis titled \textbf{``A1M: A 64-bit RISC Architecture with Complete Compiler Toolchain''} is the result of my own research and has not been submitted elsewhere for any degree or diploma.

All sources of information have been duly acknowledged.

\vspace{2cm}

\noindent
\begin{tabular}{@{}ll}
\textbf{Shahrear Hossain Shawon} & Date: \rule{3cm}{0.4pt} \\[2cm]
\textbf{Supervisor:} & \\
\textbf{Lecturer Riazul Islam} & Date: \rule{3cm}{0.4pt} \\
Department of Electrical and Electronics Engineering & \\
International Islamic University Chittagong & \\
\end{tabular}

% ============================================================================
% ABSTRACT
% ============================================================================
\chapter*{Abstract}
\addcontentsline{toc}{chapter}{Abstract}

This thesis presents A1M, a comprehensive 64-bit RISC (Reduced Instruction Set Computer) architecture with a complete compiler toolchain. The project encompasses the entire computer science stack, from hardware design in Verilog HDL to a high-level programming language with full compilation support.

The A1M architecture features a 32-register design, 5-stage pipeline, and over 60 instructions supporting arithmetic, logical, memory, and control flow operations. The hardware implementation is fully specified in Verilog and can be synthesized for FPGA deployment.

A complete software toolchain accompanies the hardware, including:
\begin{itemize}
    \item A two-pass assembler for A1M assembly language
    \item A cycle-accurate emulator for testing and debugging
    \item The Algo programming language, a C-like high-level language
    \item A complete compiler with lexer, parser, semantic analyzer, and code generator
    \item VS Code integration with syntax highlighting and debugging support
\end{itemize}

This work serves as both a functional computer system and an educational resource, demonstrating fundamental concepts in computer architecture, assembly language programming, and compiler construction. The entire system is open-source and designed for educational purposes.

\textbf{Keywords:} RISC Architecture, Compiler Design, HDL, Assembly Language, Computer Architecture, Educational Computing

\vspace{1cm}

\noindent
\textbf{Supervisor:} Lecturer Riazul Islam\\
Department of Electrical and Electronics Engineering\\
International Islamic University Chittagong

% ============================================================================
% ACKNOWLEDGMENTS
% ============================================================================
\chapter*{Acknowledgments}
\addcontentsline{toc}{chapter}{Acknowledgments}

I would like to express my sincere gratitude to all those who supported and contributed to this research project.

First and foremost, I thank \textbf{Lecturer Riazul Islam}, my supervisor, for his invaluable guidance, continuous support, and expert advice throughout the development of this work. His insights into computer architecture and compiler design were instrumental in shaping this research.

I am grateful to the Department of Electrical and Electronics Engineering at International Islamic University Chittagong for providing the research environment and resources necessary for this work.

I acknowledge the faculty members in the Department of Electrical and Electronics Engineering for their valuable feedback and encouragement throughout this project.

Special thanks to the open-source community for providing excellent tools and documentation that aided my research, particularly the developers of GCC, LLVM, and various HDL synthesis tools.

I acknowledge the support of my family and friends who encouraged me throughout this journey.

Finally, I dedicate this work to all students and researchers interested in understanding the fundamental principles of computer systems, from transistors to compilers.

\vspace{1cm}

\noindent
Shahrear Hossain Shawon\\
Department of Electrical and Electronics Engineering\\
International Islamic University Chittagong

% ============================================================================
% TABLE OF CONTENTS
% ============================================================================
\setcounter{tocdepth}{3}  % Show subsections in TOC
\setcounter{secnumdepth}{3}  % Number subsections
\tableofcontents
\newpage

\listoffigures
\newpage

\listoftables
\newpage

\lstlistoflistings
\newpage

% ============================================================================
% PART I: HARDWARE ARCHITECTURE
% ============================================================================
\part{Hardware Architecture}

% ============================================================================
% CHAPTER 1: INTRODUCTION
% ============================================================================
\chapter{Introduction to A1M Architecture}

\section{Overview}

The A1M architecture represents a complete educational computer system designed to demonstrate fundamental principles of computer architecture, assembly language programming, and compiler design. This 64-bit RISC architecture provides a clean, understandable foundation for learning how computers work at every level of abstraction.

\subsection{Design Philosophy}

The A1M architecture adheres to RISC (Reduced Instruction Set Computer) principles:

\begin{itemize}
    \item \textbf{Simplicity:} Fixed-length 32-bit instructions
    \item \textbf{Regularity:} Consistent instruction formats
    \item \textbf{Orthogonality:} Independent instruction components
    \item \textbf{Performance:} Pipeline-friendly design
    \item \textbf{Extensibility:} Room for future enhancements
\end{itemize}

\subsection{Key Features}

\begin{itemize}
    \item 64-bit data path and registers
    \item 32 general-purpose registers (R0-R31)
    \item 32-bit fixed-length instruction encoding
    \item Harvard architecture (separate instruction and data memory)
    \item 5-stage pipeline (IF, ID, EX, MEM, WB)
    \item Over 60 instructions covering all essential operations
    \item Support for 64-bit addressing (16 exabytes theoretical)
\end{itemize}

\subsection{RISC vs. CISC}

Table \ref{tab:risc-cisc} compares RISC and CISC architectures:

\begin{table}[H]
\centering
\caption{RISC vs. CISC Comparison}
\label{tab:risc-cisc}
\begin{tabularx}{\textwidth}{|X|X|}
\hline
\textbf{RISC (A1M)} & \textbf{CISC (x86)} \\
\hline
Fixed-length instructions (32-bit) & Variable-length instructions (1-15 bytes) \\
\hline
Simple addressing modes & Complex addressing modes \\
\hline
Load/Store architecture & Register-memory operations \\
\hline
Large register file (32 registers) & Small register file (8-16 registers) \\
\hline
Hardwired control & Microcode control \\
\hline
Pipeline-friendly & Pipeline stalls common \\
\hline
\end{tabularx}
\end{table}

\section{System Architecture}

Figure \ref{fig:system-overview} shows the high-level system architecture:

\begin{figure}[H]
\centering
\begin{verbatim}
┌─────────────────────────────────────────────────────────┐
│                    A1M CPU System                       │
├─────────────────────────────────────────────────────────┤
│                                                         │
│  ┌───────────┐    ┌──────────┐    ┌──────────────┐      │
│  │           │    │          │    │              │      │
│  │ Registers │◄───┤   ALU    │◄───┤ Control Unit │      │
│  │  (R0-R31) │    │          │    │              │      │
│  │           │    └──────────┘    └──────────────┘      │
│  └───────────┘                                          │
│       ▲                                                 │
│       │                                                 │
│       ▼                                                 │
│  ┌─────────────────────────────────────────────────┐    │
│  │              Memory Subsystem                   │    │
│  │  ┌─────────────────┐  ┌────────────────────┐    │    │
│  │  │ Instruction Mem │  │    Data Memory     │    │    │
│  │  │   (I-Cache)     │  │    (D-Cache)       │    │    │
│  │  └─────────────────┘  └────────────────────┘    │    │
│  └─────────────────────────────────────────────────┘    │
│                                                         │
└─────────────────────────────────────────────────────────┘
\end{verbatim}
\caption{A1M System Overview}
\label{fig:system-overview}
\end{figure}

\section{Thesis Organization}

This thesis is organized into eight parts:

\begin{description}
    \item[Part I:] Hardware Architecture --- Registers, ISA, ALU, Control Unit, Memory, Pipeline
    \item[Part II:] Assembly Language --- Assembly fundamentals, data movement, arithmetic, control flow, system calls, assembler implementation
    \item[Part III:] Emulator --- Emulator architecture and system call implementation
    \item[Part IV:] Algo Language --- High-level language introduction, expressions, control structures, functions, arrays
    \item[Part V:] Compiler Implementation --- Lexical analysis, parsing, semantic analysis, code generation
    \item[Part VI:] Compilation Process --- Build system, optimization techniques
    \item[Part VII:] Practical Examples --- Complete programs, debugging, performance analysis
    \item[Part VIII:] Testing and Verification --- Test suite, future enhancements
\end{description}

% ============================================================================
% CHAPTER 2: REGISTER ARCHITECTURE
% ============================================================================
\chapter{Register Architecture}

\section{Register File Organization}

The A1M architecture implements a 32-register file, providing ample storage for temporary values and reducing memory traffic. Each register is 64 bits wide, supporting full 64-bit integer operations.

\subsection{Register Naming Convention}

Table \ref{tab:registers} shows the complete register set:

\begin{table}[H]
\centering
\caption{A1M Register Set}
\label{tab:registers}
\begin{tabularx}{\textwidth}{|l|l|X|}
\hline
\textbf{Register} & \textbf{ABI Name} & \textbf{Purpose} \\
\hline
R0 & zero & Hardwired to 0 \\
R1 & ret/arg1 & Return value / First syscall argument \\
R2 & arg2 & Second syscall argument \\
R3 & sp & Stack pointer \\
R4-R11 & a0-a7 & Function arguments \\
R12-R30 & t0-t18 & Temporary registers \\
R31 & ra & Return address \\
\hline
\end{tabularx}
\end{table}

\subsection{Special Registers}

\begin{itemize}
    \item \textbf{R0 (zero):} Always contains the value 0. Writes to R0 are ignored. This simplifies many operations and provides a source of constant zero.
    \item \textbf{R3 (sp):} Stack pointer, points to the top of the stack. Grows downward (toward lower addresses).
    \item \textbf{R31 (ra):} Return address, automatically set by the CALL instruction. Contains the address to return to after a function completes.
\end{itemize}

\section{Verilog Implementation}

The register file is implemented in Verilog with dual read ports and one write port:

\begin{lstlisting}[style=verilog, caption={Register File Implementation}, label={lst:regfile}]
module register_file (
    input wire clk,
    input wire reset,
    
    // Read ports
    input wire [4:0] read_addr1,
    input wire [4:0] read_addr2,
    output reg [63:0] read_data1,
    output reg [63:0] read_data2,
    
    // Write port
    input wire write_enable,
    input wire [4:0] write_addr,
    input wire [63:0] write_data
);

// Register array
reg [63:0] registers [0:31];

integer i;

// Initialize registers
initial begin
    for (i = 0; i < 32; i = i + 1) begin
        registers[i] = 64'h0;
    end
end

// Asynchronous read (combinational)
always @(*) begin
    if (read_addr1 == 5'h0)
        read_data1 = 64'h0;  // R0 always zero
    else
        read_data1 = registers[read_addr1];
    
    if (read_addr2 == 5'h0)
        read_data2 = 64'h0;  // R0 always zero
    else
        read_data2 = registers[read_addr2];
end

// Synchronous write (on rising edge)
always @(posedge clk) begin
    if (reset) begin
        for (i = 0; i < 32; i = i + 1) begin
            registers[i] <= 64'h0;
        end
    end else if (write_enable && write_addr != 5'h0) begin
        // Don't write to R0
        registers[write_addr] <= write_data;
    end
end

endmodule
\end{lstlisting}

\section{Register Usage Convention}

\subsection{Caller-Saved vs. Callee-Saved}

\begin{itemize}
    \item \textbf{Caller-saved:} R1-R2, R4-R11, R12-R30, R31 --- The calling function must save these if their values are needed after a function call
    \item \textbf{Callee-saved:} R3 --- Called functions must preserve these registers
\end{itemize}

\subsection{Function Call Example}

\begin{lstlisting}[style=assembly, caption={Register Usage in Function Calls}]
; Caller
caller:
    PUSH R31            ; Save return address
    PUSH R4             ; Save argument registers if needed
    
    MOVI R4, 10         ; First argument
    MOVI R5, 20         ; Second argument
    CALL add_function   ; Call function
    
    MOV R10, R1         ; Save return value
    POP R4              ; Restore saved registers
    POP R31
    RET

; Callee
add_function:
    PUSH R31            ; Save return address
    ADD R1, R4, R5      ; Compute result (R1 = R4 + R5)
    POP R31             ; Restore return address
    RET
\end{lstlisting}

% Continue with Chapter 3...

% ============================================================================
% CHAPTER 3: INSTRUCTION SET ARCHITECTURE
% ============================================================================
\chapter{Instruction Set Architecture}

\section{Instruction Formats}

The A1M ISA uses 32-bit fixed-length instructions in four formats:

\subsection{R-Type (Register)}

Used for register-to-register operations:

\begin{verbatim}
 31    24 23  19 18  14 13   9 8       0
┌────────┬──────┬──────┬──────┬─────────┐
│ Opcode │  RD  │ RS1  │ RS2  │ Unused  │
└────────┴──────┴──────┴──────┴─────────┘
  8 bits  5 bits 5 bits 5 bits  9 bits
\end{verbatim}

\begin{itemize}
    \item \textbf{Opcode:} 8-bit operation code
    \item \textbf{RD:} 5-bit destination register
    \item \textbf{RS1:} 5-bit source register 1
    \item \textbf{RS2:} 5-bit source register 2
\end{itemize}

\subsection{I-Type (Immediate)}

Used for operations with immediate values:

\begin{verbatim}
 31    24 23  19 18  14 13           0
┌────────┬──────┬──────┬──────────────┐
│ Opcode │  RD  │ RS1  │  Immediate   │
└────────┴──────┴──────┴──────────────┘
  8 bits  5 bits 5 bits   14 bits
\end{verbatim}

\begin{itemize}
    \item \textbf{Immediate:} 14-bit signed immediate value (sign-extended to 64 bits)
    \item Range: $-8192$ to $+8191$
\end{itemize}

\subsection{J-Type (Jump)}

Used for jump and branch instructions:

\begin{verbatim}
 31    24 23                      0
┌────────┬─────────────────────────┐
│ Opcode │    Jump Address         │
└────────┴─────────────────────────┘
  8 bits        24 bits
\end{verbatim}

\begin{itemize}
    \item \textbf{Jump Address:} 24-bit target address
    \item Address range: $0$ to $16,777,215$ (16 MB)
\end{itemize}

\subsection{M-Type (Memory)}

Used for memory operations (overlaps with I-Type):

\begin{verbatim}
 31    24 23  19 18  14 13           0
┌────────┬──────┬──────┬──────────────┐
│ Opcode │  RD  │ Base │    Offset    │
└────────┴──────┴──────┴──────────────┘
  8 bits  5 bits 5 bits   14 bits
\end{verbatim}

\section{Opcode Assignments}

Table \ref{tab:opcodes} shows the complete opcode assignment:

\begin{table}[H]
\centering
\caption{A1M Opcode Assignment}
\label{tab:opcodes}
\small
\begin{tabularx}{\textwidth}{|l|l|l|X|}
\hline
\textbf{Code} & \textbf{Mnemonic} & \textbf{Format} & \textbf{Operation} \\
\hline
0x00 & NOP & - & No operation \\
0x01 & HALT & - & Halt execution \\
\hline
0x02 & ADD & R & RD = RS1 + RS2 \\
0x03 & ADDI & I & RD = RS1 + IMM \\
0x04 & SUB & R & RD = RS1 - RS2 \\
0x05 & SUBI & I & RD = RS1 - IMM \\
0x06 & MUL & R & RD = RS1 $\times$ RS2 \\
0x07 & MULI & I & RD = RS1 $\times$ IMM \\
0x08 & DIV & R & RD = RS1 / RS2 \\
0x09 & DIVI & I & RD = RS1 / IMM \\
0x0A & MOD & R & RD = RS1 \% RS2 \\
0x0B & MODI & I & RD = RS1 \% IMM \\
\hline
0x0C & AND & R & RD = RS1 \& RS2 \\
0x0D & ANDI & I & RD = RS1 \& IMM \\
0x0E & OR & R & RD = RS1 | RS2 \\
0x0F & ORI & I & RD = RS1 | IMM \\
0x10 & XOR & R & RD = RS1 $\oplus$ RS2 \\
0x11 & XORI & I & RD = RS1 $\oplus$ IMM \\
0x12 & SHL & R & RD = RS1 $<<$ RS2 \\
0x13 & SHLI & I & RD = RS1 $<<$ IMM \\
0x14 & SHR & R & RD = RS1 $>>$ RS2 (logical) \\
0x15 & SHRI & I & RD = RS1 $>>$ IMM (logical) \\
0x16 & SAR & R & RD = RS1 $>>$ RS2 (arithmetic) \\
0x17 & SARI & I & RD = RS1 $>>$ IMM (arithmetic) \\
\hline
0x20 & LOAD & M & RD = Memory[RS1 + Offset] \\
0x21 & STORE & M & Memory[RS1 + Offset] = RD \\
0x22 & LOADB & M & RD = (byte)Memory[RS1 + Offset] \\
0x23 & STOREB & M & Memory[RS1 + Offset] = (byte)RD \\
\hline
0x30 & JMP & J & PC = Address \\
0x31 & JZ & J & if (ZF) PC = Address \\
0x32 & JNZ & J & if (!ZF) PC = Address \\
0x33 & JGT & J & if (!ZF \&\& !NF) PC = Address \\
0x34 & JLT & J & if (NF) PC = Address \\
0x39 & CALL & J & R31 = PC + 4; PC = Address \\
0x3A & RET & - & PC = R31 \\
\hline
0x40 & CMP & R & Set flags: RS1 - RS2 \\
0x41 & CMPI & I & Set flags: RS1 - IMM \\
0x42 & PUSH & R & SP -= 8; Memory[SP] = RS \\
0x43 & POP & R & RD = Memory[SP]; SP += 8 \\
\hline
0x50 & SYSCALL & I & System call (IMM = number) \\
\hline
\end{tabularx}
\end{table}

\section{Instruction Encoding Examples}

\subsection{Example 1: ADD R1, R2, R3}

\begin{verbatim}
Instruction: ADD R1, R2, R3  (R1 = R2 + R3)
Format: R-Type

Binary encoding:
 Opcode    RD     RS1    RS2   Unused
 00000010  00001  00010  00011 000000000
 
Hexadecimal: 0x02 11 23 00 = 0x02112300
\end{verbatim}

\subsection{Example 2: ADDI R4, R5, 100}

\begin{verbatim}
Instruction: ADDI R4, R5, 100  (R4 = R5 + 100)
Format: I-Type

Binary encoding:
 Opcode    RD     RS1    Immediate
 00000011  00100  00101  00000001100100
 
Hexadecimal: 0x03 14 28 64 = 0x03142864
\end{verbatim}

% ============================================================================
% CHAPTER 4: ARITHMETIC LOGIC UNIT (ALU)
% ============================================================================
\chapter{Arithmetic Logic Unit (ALU)}

\section{ALU Overview}

The Arithmetic Logic Unit (ALU) is the computational heart of the A1M processor. It performs all arithmetic and logical operations on 64-bit operands.

\subsection{Supported Operations}

The A1M ALU supports the following operations:

\begin{itemize}
    \item \textbf{Arithmetic:} ADD, SUB, MUL, DIV, MOD
    \item \textbf{Logical:} AND, OR, XOR, NOT
    \item \textbf{Shift:} SHL (left), SHR (logical right), SAR (arithmetic right)
    \item \textbf{Comparison:} CMP (sets flags)
\end{itemize}

\subsection{ALU Control Signals}

Table \ref{tab:alu-ops} shows the ALU operation encoding:

\begin{table}[H]
\centering
\caption{ALU Operation Codes}
\label{tab:alu-ops}
\begin{tabularx}{\textwidth}{|l|l|X|}
\hline
\textbf{Code} & \textbf{Operation} & \textbf{Description} \\
\hline
4'b0000 & ADD & Addition: result = A + B \\
4'b0001 & SUB & Subtraction: result = A - B \\
4'b0010 & MUL & Multiplication: result = A $\times$ B \\
4'b0011 & DIV & Division: result = A / B \\
4'b0100 & MOD & Modulo: result = A \% B \\
4'b0101 & AND & Bitwise AND: result = A \& B \\
4'b0110 & OR & Bitwise OR: result = A | B \\
4'b0111 & XOR & Bitwise XOR: result = A $\oplus$ B \\
4'b1000 & SHL & Shift left: result = A $<<$ B \\
4'b1001 & SHR & Shift right (logical): result = A $>>$ B \\
4'b1010 & SAR & Shift right (arithmetic) \\
4'b1011 & CMP & Compare: set flags based on A - B \\
\hline
\end{tabularx}
\end{table}

\section{Verilog Implementation}

\begin{lstlisting}[style=verilog, caption={ALU Implementation}, label={lst:alu}]
module alu (
    input wire [63:0] operand_a,
    input wire [63:0] operand_b,
    input wire [3:0] alu_op,
    output reg [63:0] result,
    output reg zero_flag,
    output reg negative_flag,
    output reg carry_flag,
    output reg overflow_flag
);

// ALU operation codes
localparam ALU_ADD = 4'b0000;
localparam ALU_SUB = 4'b0001;
localparam ALU_MUL = 4'b0010;
localparam ALU_DIV = 4'b0011;
localparam ALU_MOD = 4'b0100;
localparam ALU_AND = 4'b0101;
localparam ALU_OR  = 4'b0110;
localparam ALU_XOR = 4'b0111;
localparam ALU_SHL = 4'b1000;
localparam ALU_SHR = 4'b1001;
localparam ALU_SAR = 4'b1010;
localparam ALU_CMP = 4'b1011;

reg [64:0] temp_result; // For carry detection

always @(*) begin
    // Default flag values
    carry_flag = 1'b0;
    overflow_flag = 1'b0;
    
    case (alu_op)
        ALU_ADD: begin
            temp_result = {1'b0, operand_a} + {1'b0, operand_b};
            result = temp_result[63:0];
            carry_flag = temp_result[64];
            overflow_flag = (operand_a[63] == operand_b[63]) && 
                          (result[63] != operand_a[63]);
        end
        
        ALU_SUB: begin
            temp_result = {1'b0, operand_a} - {1'b0, operand_b};
            result = temp_result[63:0];
            carry_flag = temp_result[64];
            overflow_flag = (operand_a[63] != operand_b[63]) && 
                          (result[63] != operand_a[63]);
        end
        
        ALU_MUL: begin
            result = operand_a * operand_b;
        end
        
        ALU_DIV: begin
            result = (operand_b != 0) ? (operand_a / operand_b) : 64'h0;
        end
        
        ALU_MOD: begin
            result = (operand_b != 0) ? (operand_a % operand_b) : 64'h0;
        end
        
        ALU_AND: begin
            result = operand_a & operand_b;
        end
        
        ALU_OR: begin
            result = operand_a | operand_b;
        end
        
        ALU_XOR: begin
            result = operand_a ^ operand_b;
        end
        
        ALU_SHL: begin
            result = operand_a << operand_b[5:0]; // Limit shift to 0-63
        end
        
        ALU_SHR: begin
            result = operand_a >> operand_b[5:0];
        end
        
        ALU_SAR: begin
            result = $signed(operand_a) >>> operand_b[5:0];
        end
        
        ALU_CMP: begin
            temp_result = {1'b0, operand_a} - {1'b0, operand_b};
            result = temp_result[63:0];
        end
        
        default: begin
            result = 64'h0;
        end
    endcase
    
    // Update flags
    zero_flag = (result == 64'h0);
    negative_flag = result[63];
end

endmodule
\end{lstlisting}

\section{Flag Register}

The ALU sets four condition flags after each operation:

\begin{description}
    \item[Zero Flag (ZF):] Set when result is zero
    \item[Negative Flag (NF):] Set when result is negative (bit 63 = 1)
    \item[Carry Flag (CF):] Set on unsigned overflow
    \item[Overflow Flag (OF):] Set on signed overflow
\end{description}

These flags are used by conditional branch instructions to make decisions.

% ============================================================================
% CHAPTER 5: CONTROL UNIT
% ============================================================================
\chapter{Control Unit}

\section{Control Unit Architecture}

The control unit decodes instructions and generates control signals for the datapath. It is implemented as a hardwired control unit rather than microcode for simplicity and performance.

\section{Control Signals}

Table \ref{tab:control-signals} lists all control signals:

\begin{table}[H]
\centering
\caption{Control Signals}
\label{tab:control-signals}
\small
\begin{tabularx}{\textwidth}{|l|l|X|}
\hline
\textbf{Signal} & \textbf{Width} & \textbf{Purpose} \\
\hline
reg\_write & 1 & Enable register write \\
alu\_op & 4 & ALU operation code \\
alu\_src & 1 & ALU source: 0=reg, 1=imm \\
mem\_read & 1 & Enable memory read \\
mem\_write & 1 & Enable memory write \\
mem\_to\_reg & 1 & Write memory to register \\
branch & 1 & Branch instruction \\
jump & 1 & Unconditional jump \\
pc\_src & 2 & PC source: 00=PC+4, 01=branch, 10=jump, 11=register \\
\hline
\end{tabularx}
\end{table}

\section{Verilog Implementation}

\begin{lstlisting}[style=verilog, caption={Control Unit Implementation}, label={lst:control}]
module control_unit (
    input wire [7:0] opcode,
    output reg reg_write,
    output reg [3:0] alu_op,
    output reg alu_src,
    output reg mem_read,
    output reg mem_write,
    output reg mem_to_reg,
    output reg branch,
    output reg jump,
    output reg [1:0] pc_src
);

// Opcode definitions
localparam OP_NOP    = 8'h00;
localparam OP_HALT   = 8'h01;
localparam OP_ADD    = 8'h02;
localparam OP_ADDI   = 8'h03;
localparam OP_SUB    = 8'h04;
localparam OP_SUBI   = 8'h05;
localparam OP_MUL    = 8'h06;
localparam OP_DIV    = 8'h08;
localparam OP_AND    = 8'h0C;
localparam OP_ANDI   = 8'h0D;
localparam OP_OR     = 8'h0E;
localparam OP_XOR    = 8'h10;
localparam OP_SHL    = 8'h12;
localparam OP_SHR    = 8'h14;
localparam OP_LOAD   = 8'h20;
localparam OP_STORE  = 8'h21;
localparam OP_JMP    = 8'h30;
localparam OP_JZ     = 8'h31;
localparam OP_JNZ    = 8'h32;
localparam OP_JGT    = 8'h33;
localparam OP_JLT    = 8'h34;
localparam OP_CALL   = 8'h39;
localparam OP_RET    = 8'h3A;
localparam OP_CMP    = 8'h40;
localparam OP_PUSH   = 8'h42;
localparam OP_POP    = 8'h43;

always @(*) begin
    // Default values
    reg_write = 1'b0;
    alu_op = 4'b0000;
    alu_src = 1'b0;
    mem_read = 1'b0;
    mem_write = 1'b0;
    mem_to_reg = 1'b0;
    branch = 1'b0;
    jump = 1'b0;
    pc_src = 2'b00;
    
    case (opcode)
        OP_NOP: begin
            // No operation
        end
        
        OP_ADD: begin
            reg_write = 1'b1;
            alu_op = 4'b0000; // ADD
            alu_src = 1'b0;   // Register source
        end
        
        OP_ADDI: begin
            reg_write = 1'b1;
            alu_op = 4'b0000; // ADD
            alu_src = 1'b1;   // Immediate source
        end
        
        OP_SUB: begin
            reg_write = 1'b1;
            alu_op = 4'b0001; // SUB
            alu_src = 1'b0;
        end
        
        OP_SUBI: begin
            reg_write = 1'b1;
            alu_op = 4'b0001; // SUB
            alu_src = 1'b1;
        end
        
        OP_MUL: begin
            reg_write = 1'b1;
            alu_op = 4'b0010; // MUL
            alu_src = 1'b0;
        end
        
        OP_DIV: begin
            reg_write = 1'b1;
            alu_op = 4'b0011; // DIV
            alu_src = 1'b0;
        end
        
        OP_AND: begin
            reg_write = 1'b1;
            alu_op = 4'b0101; // AND
            alu_src = 1'b0;
        end
        
        OP_OR: begin
            reg_write = 1'b1;
            alu_op = 4'b0110; // OR
            alu_src = 1'b0;
        end
        
        OP_XOR: begin
            reg_write = 1'b1;
            alu_op = 4'b0111; // XOR
            alu_src = 1'b0;
        end
        
        OP_SHL: begin
            reg_write = 1'b1;
            alu_op = 4'b1000; // SHL
            alu_src = 1'b0;
        end
        
        OP_SHR: begin
            reg_write = 1'b1;
            alu_op = 4'b1001; // SHR
            alu_src = 1'b0;
        end
        
        OP_LOAD: begin
            reg_write = 1'b1;
            mem_read = 1'b1;
            mem_to_reg = 1'b1;
            alu_op = 4'b0000; // ADD for address calculation
            alu_src = 1'b1;   // Use immediate offset
        end
        
        OP_STORE: begin
            mem_write = 1'b1;
            alu_op = 4'b0000; // ADD for address calculation
            alu_src = 1'b1;
        end
        
        OP_JMP: begin
            jump = 1'b1;
            pc_src = 2'b10;
        end
        
        OP_JZ, OP_JNZ, OP_JGT, OP_JLT: begin
            branch = 1'b1;
            pc_src = 2'b01;
        end
        
        OP_CALL: begin
            reg_write = 1'b1; // Write return address to R31
            jump = 1'b1;
            pc_src = 2'b10;
        end
        
        OP_RET: begin
            jump = 1'b1;
            pc_src = 2'b11; // PC from R31
        end
        
        OP_CMP: begin
            alu_op = 4'b1011; // CMP
            alu_src = 1'b0;
        end
        
        OP_PUSH: begin
            mem_write = 1'b1;
            reg_write = 1'b1; // Update SP
        end
        
        OP_POP: begin
            mem_read = 1'b1;
            reg_write = 1'b1;
            mem_to_reg = 1'b1;
        end
        
        default: begin
            // Invalid opcode - all signals remain at default
        end
    endcase
end

endmodule
\end{lstlisting}

% ============================================================================
% CHAPTER 6: MEMORY SUBSYSTEM
% ============================================================================
\chapter{Memory Subsystem}

\section{Memory Architecture}

The A1M architecture uses a Harvard architecture with separate instruction and data memories. This allows simultaneous instruction fetch and data access, improving pipeline performance.

\subsection{Memory Organization}

\begin{itemize}
    \item \textbf{Address Width:} 64 bits (theoretical 16 exabytes)
    \item \textbf{Data Width:} 64 bits (8 bytes per word)
    \item \textbf{Instruction Memory:} Read-only during execution
    \item \textbf{Data Memory:} Read/Write
    \item \textbf{Byte Addressable:} Each byte has a unique address
\end{itemize}

\subsection{Memory Map}

Table \ref{tab:memory-map} shows the default memory layout:

\begin{table}[H]
\centering
\caption{A1M Memory Map}
\label{tab:memory-map}
\begin{tabularx}{\textwidth}{|l|l|X|}
\hline
\textbf{Address Range} & \textbf{Region} & \textbf{Purpose} \\
\hline
0x0000000000000000 - 0x00000000000FFFFF & Code & Instruction memory (1 MB) \\
0x0000000000100000 - 0x00000000001FFFFF & Data & Global/static data (1 MB) \\
0x0000000000200000 - 0x00000000002FFFFF & Heap & Dynamic allocation (1 MB) \\
0x00000000FFFF0000 - 0x00000000FFFFFFFF & Stack & Stack (grows down, 64 KB) \\
\hline
\end{tabularx}
\end{table}

\section{Verilog Implementation}

\begin{lstlisting}[style=verilog, caption={Data Memory Module}, label={lst:data-mem}]
module data_memory #(
    parameter ADDR_WIDTH = 16,  // 64K words = 512KB
    parameter DATA_WIDTH = 64
)(
    input wire clk,
    input wire mem_read,
    input wire mem_write,
    input wire [ADDR_WIDTH-1:0] address,
    input wire [DATA_WIDTH-1:0] write_data,
    output reg [DATA_WIDTH-1:0] read_data
);

// Memory array (64K 64-bit words = 512KB)
reg [DATA_WIDTH-1:0] memory [0:(1 << ADDR_WIDTH)-1];

// Initialize memory to zero
integer i;
initial begin
    for (i = 0; i < (1 << ADDR_WIDTH); i = i + 1) begin
        memory[i] = 64'h0;
    end
end

// Synchronous write, asynchronous read
always @(posedge clk) begin
    if (mem_write) begin
        memory[address] <= write_data;
    end
end

always @(*) begin
    if (mem_read) begin
        read_data = memory[address];
    end else begin
        read_data = 64'h0;
    end
end

endmodule
\end{lstlisting}

\begin{lstlisting}[style=verilog, caption={Instruction Memory Module}, label={lst:inst-mem}]
module instruction_memory #(
    parameter ADDR_WIDTH = 16,
    parameter DATA_WIDTH = 32
)(
    input wire [ADDR_WIDTH-1:0] address,
    output reg [DATA_WIDTH-1:0] instruction
);

// Instruction memory array
reg [DATA_WIDTH-1:0] memory [0:(1 << ADDR_WIDTH)-1];

// Load program from file (simulation)
initial begin
    $readmemh("program.hex", memory);
end

// Asynchronous read
always @(*) begin
    instruction = memory[address];
end

endmodule
\end{lstlisting}

\section{Cache Design (Future Enhancement)}

The current implementation uses simple memory modules. A production design would include:

\begin{itemize}
    \item \textbf{L1 I-Cache:} 32 KB, 4-way set associative
    \item \textbf{L1 D-Cache:} 32 KB, 4-way set associative
    \item \textbf{L2 Unified Cache:} 256 KB, 8-way set associative
    \item \textbf{Cache Line Size:} 64 bytes
    \item \textbf{Write Policy:} Write-back with write-allocate
\end{itemize}

% ============================================================================
% CHAPTER 7: PIPELINE ARCHITECTURE
% ============================================================================
\chapter{Pipeline Architecture}

\section{Five-Stage Pipeline}

The A1M processor implements a classic 5-stage RISC pipeline:

\begin{enumerate}
    \item \textbf{IF (Instruction Fetch):} Fetch instruction from memory
    \item \textbf{ID (Instruction Decode):} Decode instruction and read registers
    \item \textbf{EX (Execute):} Execute ALU operation
    \item \textbf{MEM (Memory):} Access data memory
    \item \textbf{WB (Write Back):} Write result to register
\end{enumerate}

\subsection{Pipeline Diagram}

\begin{verbatim}
Time ->  1    2    3    4    5    6    7
Inst 1: IF   ID   EX  MEM   WB
Inst 2:      IF   ID   EX  MEM   WB
Inst 3:           IF   ID   EX  MEM   WB
Inst 4:                IF   ID   EX  MEM   WB
\end{verbatim}

In ideal conditions, the pipeline achieves a CPI (Cycles Per Instruction) of 1.0 after the initial latency.

\section{Pipeline Hazards}

\subsection{Data Hazards}

Data hazards occur when an instruction depends on the result of a previous instruction still in the pipeline.

\textbf{Example:}
\begin{lstlisting}[style=assembly]
ADD R1, R2, R3    ; R1 = R2 + R3
SUB R4, R1, R5    ; Needs R1 (not yet written!)
\end{lstlisting}

\textbf{Solutions:}
\begin{itemize}
    \item \textbf{Forwarding:} Pass result directly from EX/MEM to EX stage
    \item \textbf{Stalling:} Insert NOP bubbles when forwarding is insufficient
\end{itemize}

\subsection{Control Hazards}

Control hazards occur with branch and jump instructions.

\textbf{Solutions:}
\begin{itemize}
    \item \textbf{Branch Prediction:} Predict branch not taken (simple)
    \item \textbf{Branch Delay Slot:} Execute next instruction regardless
    \item \textbf{Pipeline Flush:} Cancel instructions if prediction wrong
\end{itemize}

\section{Forwarding Unit}

\begin{lstlisting}[style=verilog, caption={Forwarding Unit}, label={lst:forward}]
module forwarding_unit (
    // EX/MEM pipeline register
    input wire [4:0] ex_mem_rd,
    input wire ex_mem_reg_write,
    
    // MEM/WB pipeline register
    input wire [4:0] mem_wb_rd,
    input wire mem_wb_reg_write,
    
    // ID/EX pipeline register
    input wire [4:0] id_ex_rs1,
    input wire [4:0] id_ex_rs2,
    
    // Forwarding control
    output reg [1:0] forward_a,  // 00=no forward, 01=EX/MEM, 10=MEM/WB
    output reg [1:0] forward_b
);

always @(*) begin
    // Default: no forwarding
    forward_a = 2'b00;
    forward_b = 2'b00;
    
    // EX hazard (forward from EX/MEM)
    if (ex_mem_reg_write && (ex_mem_rd != 5'h0)) begin
        if (ex_mem_rd == id_ex_rs1)
            forward_a = 2'b01;
        if (ex_mem_rd == id_ex_rs2)
            forward_b = 2'b01;
    end
    
    // MEM hazard (forward from MEM/WB)
    if (mem_wb_reg_write && (mem_wb_rd != 5'h0)) begin
        if ((mem_wb_rd == id_ex_rs1) && 
            !(ex_mem_reg_write && (ex_mem_rd == id_ex_rs1)))
            forward_a = 2'b10;
        if ((mem_wb_rd == id_ex_rs2) && 
            !(ex_mem_reg_write && (ex_mem_rd == id_ex_rs2)))
            forward_b = 2'b10;
    end
end

endmodule
\end{lstlisting}

\section{Hazard Detection Unit}

\begin{lstlisting}[style=verilog, caption={Hazard Detection Unit}, label={lst:hazard}]
module hazard_detection (
    // ID/EX pipeline register
    input wire id_ex_mem_read,
    input wire [4:0] id_ex_rd,
    
    // IF/ID pipeline register
    input wire [4:0] if_id_rs1,
    input wire [4:0] if_id_rs2,
    
    // Control signals
    output reg stall,
    output reg pc_write,
    output reg if_id_write
);

always @(*) begin
    // Default: no stall
    stall = 1'b0;
    pc_write = 1'b1;
    if_id_write = 1'b1;
    
    // Load-use hazard detection
    if (id_ex_mem_read && 
        ((id_ex_rd == if_id_rs1) || (id_ex_rd == if_id_rs2))) begin
        // Stall the pipeline
        stall = 1'b1;
        pc_write = 1'b0;      // Don't update PC
        if_id_write = 1'b0;   // Don't update IF/ID
    end
end

endmodule
\end{lstlisting}

% ============================================================================
% PART II: ASSEMBLY LANGUAGE
% ============================================================================
\part{Assembly Language Programming}

% ============================================================================
% CHAPTER 8: ASSEMBLY LANGUAGE FUNDAMENTALS
% ============================================================================
\chapter{Assembly Language Fundamentals}

\section{Introduction to A1M Assembly}

Assembly language provides a human-readable representation of machine code. Each assembly instruction corresponds directly to a machine instruction.

\subsection{Assembly Syntax}

The A1M assembly language follows this general syntax:

\begin{verbatim}
label:  INSTRUCTION  operands    ; comment
\end{verbatim}

\textbf{Example:}
\begin{lstlisting}[style=assembly, caption={Simple Assembly Program}]
; Hello World program
.section .data
msg:    .string "Hello, A1M!\n"
len:    .quad 12

.section .text
.global _start

_start:
    ; Write system call
    MOVI R1, 1          ; syscall number (write)
    MOVI R2, 1          ; file descriptor (stdout)
    LOAD R3, msg        ; buffer address
    LOAD R4, len        ; message length
    SYSCALL 0           ; invoke system call
    
    ; Exit system call
    MOVI R1, 0          ; syscall number (exit)
    MOVI R2, 0          ; exit code
    SYSCALL 0
\end{lstlisting}

\section{Directives}

Assembly directives provide instructions to the assembler (not the CPU):

\begin{table}[H]
\centering
\caption{Assembly Directives}
\label{tab:directives}
\begin{tabularx}{\textwidth}{|l|X|}
\hline
\textbf{Directive} & \textbf{Purpose} \\
\hline
.section .text & Begin code section \\
.section .data & Begin data section \\
.section .bss & Begin uninitialized data section \\
.global label & Make label globally visible \\
.byte value & Define byte (8 bits) \\
.word value & Define word (32 bits) \\
.quad value & Define quadword (64 bits) \\
.string "text" & Define null-terminated string \\
.space n & Reserve n bytes of space \\
.align n & Align to n-byte boundary \\
\hline
\end{tabularx}
\end{table}

\section{Labels and Symbols}

Labels mark positions in code or data:

\begin{lstlisting}[style=assembly]
loop_start:
    ADDI R10, R10, 1
    CMPI R10, 100
    JLT loop_start      ; Jump to label
\end{lstlisting}

% ============================================================================
% CHAPTER 9: DATA MOVEMENT INSTRUCTIONS
% ============================================================================
\chapter{Data Movement Instructions}

\section{Register Operations}

\subsection{MOV - Move Register}

Copy value from one register to another:

\begin{lstlisting}[style=assembly]
MOV R1, R2          ; R1 = R2
\end{lstlisting}

Implemented as: \texttt{ADD R1, R2, R0}

\subsection{MOVI - Move Immediate}

Load immediate value into register:

\begin{lstlisting}[style=assembly]
MOVI R1, 42         ; R1 = 42
MOVI R2, -100       ; R2 = -100
\end{lstlisting}

\section{Memory Operations}

\subsection{LOAD - Load from Memory}

Load 64-bit value from memory:

\begin{lstlisting}[style=assembly]
LOAD R1, 0(R3)      ; R1 = Memory[R3 + 0]
LOAD R2, 16(R3)     ; R2 = Memory[R3 + 16]
LOAD R4, array      ; R4 = Memory[array address]
\end{lstlisting}

\subsection{STORE - Store to Memory}

Store 64-bit value to memory:

\begin{lstlisting}[style=assembly]
STORE R1, 0(R3)     ; Memory[R3 + 0] = R1
STORE R2, 16(R3)    ; Memory[R3 + 16] = R2
\end{lstlisting}

\subsection{LOADB/STOREB - Byte Operations}

Load or store single byte:

\begin{lstlisting}[style=assembly]
LOADB R1, 0(R3)     ; R1 = (byte)Memory[R3]
STOREB R2, 0(R4)    ; Memory[R4] = (byte)R2
\end{lstlisting}

\section{Stack Operations}

\subsection{PUSH - Push to Stack}

\begin{lstlisting}[style=assembly]
PUSH R1             ; SP -= 8; Memory[SP] = R1
\end{lstlisting}

Equivalent to:
\begin{lstlisting}[style=assembly]
SUBI R3, R3, 8      ; R3 (SP) -= 8
STORE R1, 0(R3)     ; Memory[SP] = R1
\end{lstlisting}

\subsection{POP - Pop from Stack}

\begin{lstlisting}[style=assembly]
POP R1              ; R1 = Memory[SP]; SP += 8
\end{lstlisting}

\section{Example: Array Access}

\begin{lstlisting}[style=assembly, caption={Array Access Example}]
.section .data
array:  .quad 10, 20, 30, 40, 50

.section .text
_start:
    ; Load base address
    MOVI R10, array     ; R10 = &array[0]
    
    ; Access array[2] (third element)
    MOVI R11, 2         ; Index
    MOVI R12, 8         ; Size of quad (8 bytes)
    MUL R11, R11, R12   ; Offset = index * size
    ADD R13, R10, R11   ; Address = base + offset
    LOAD R14, 0(R13)    ; R14 = array[2] = 30
\end{lstlisting}

% ============================================================================
% CHAPTER 10: ARITHMETIC OPERATIONS
% ============================================================================
\chapter{Arithmetic Operations}

\section{Basic Arithmetic}

\subsection{Addition}

\begin{lstlisting}[style=assembly]
ADD R1, R2, R3      ; R1 = R2 + R3 (register)
ADDI R1, R2, 100    ; R1 = R2 + 100 (immediate)
\end{lstlisting}

\subsection{Subtraction}

\begin{lstlisting}[style=assembly]
SUB R1, R2, R3      ; R1 = R2 - R3
SUBI R1, R2, 50     ; R1 = R2 - 50
\end{lstlisting}

\subsection{Multiplication}

\begin{lstlisting}[style=assembly]
MUL R1, R2, R3      ; R1 = R2 * R3
MULI R1, R2, 10     ; R1 = R2 * 10
\end{lstlisting}

\subsection{Division and Modulo}

\begin{lstlisting}[style=assembly]
DIV R1, R2, R3      ; R1 = R2 / R3 (quotient)
MOD R1, R2, R3      ; R1 = R2 % R3 (remainder)

DIVI R1, R10, 10    ; R1 = R10 / 10
MODI R2, R10, 10    ; R2 = R10 % 10
\end{lstlisting}

\section{Logical Operations}

\subsection{Bitwise AND, OR, XOR}

\begin{lstlisting}[style=assembly]
AND R1, R2, R3      ; R1 = R2 & R3
ANDI R1, R2, 0xFF   ; R1 = R2 & 0xFF (mask lower 8 bits)

OR R1, R2, R3       ; R1 = R2 | R3
ORI R1, R2, 0x80    ; R1 = R2 | 0x80 (set bit 7)

XOR R1, R2, R3      ; R1 = R2 ^ R3
XORI R1, R1, -1     ; R1 = ~R1 (NOT operation)
\end{lstlisting}

\subsection{Shift Operations}

\begin{lstlisting}[style=assembly]
SHL R1, R2, R3      ; R1 = R2 << R3 (left shift)
SHLI R1, R2, 3      ; R1 = R2 << 3

SHR R1, R2, R3      ; R1 = R2 >> R3 (logical right)
SHRI R1, R2, 2      ; R1 = R2 >> 2

SAR R1, R2, R3      ; R1 = R2 >> R3 (arithmetic right)
SARI R1, R2, 4      ; R1 = R2 >> 4 (preserve sign)
\end{lstlisting}

\section{Example: Calculate Expression}

Calculate: $result = (a + b) \times c - d / 2$

\begin{lstlisting}[style=assembly, caption={Complex Expression}]
; Assume: R10=a, R11=b, R12=c, R13=d

; Step 1: temp1 = a + b
ADD R14, R10, R11       ; R14 = a + b

; Step 2: temp2 = temp1 * c
MUL R15, R14, R12       ; R15 = (a + b) * c

; Step 3: temp3 = d / 2
DIVI R16, R13, 2        ; R16 = d / 2

; Step 4: result = temp2 - temp3
SUB R1, R15, R16        ; R1 = (a+b)*c - d/2
\end{lstlisting}

% ============================================================================
% CHAPTER 11: CONTROL FLOW
% ============================================================================
\chapter{Control Flow Instructions}

\section{Unconditional Jumps}

\subsection{JMP - Jump}

Transfer control unconditionally:

\begin{lstlisting}[style=assembly]
JMP target          ; PC = target address
\end{lstlisting}

\subsection{CALL - Function Call}

Call a function (saves return address):

\begin{lstlisting}[style=assembly]
CALL function       ; R31 = PC + 4; PC = function
\end{lstlisting}

\subsection{RET - Return}

Return from function:

\begin{lstlisting}[style=assembly]
RET                 ; PC = R31
\end{lstlisting}

\section{Conditional Branches}

\subsection{Comparison}

Compare two values and set flags:

\begin{lstlisting}[style=assembly]
CMP R1, R2          ; Set flags based on R1 - R2
CMPI R1, 100        ; Set flags based on R1 - 100
\end{lstlisting}

\subsection{Conditional Jumps}

\begin{lstlisting}[style=assembly]
JZ label            ; Jump if Zero Flag set
JNZ label           ; Jump if Zero Flag clear
JGT label           ; Jump if Greater Than (signed)
JLT label           ; Jump if Less Than (signed)
JGE label           ; Jump if Greater or Equal
JLE label           ; Jump if Less or Equal
\end{lstlisting}

\section{Control Flow Examples}

\subsection{If-Then-Else}

\begin{lstlisting}[style=assembly, caption={If-Then-Else Structure}]
; if (R10 > 50) {
;     R1 = 1;
; } else {
;     R1 = 0;
; }

    CMPI R10, 50        ; Compare R10 with 50
    JLE else_block      ; If R10 <= 50, go to else
    
then_block:
    MOVI R1, 1          ; R1 = 1
    JMP end_if          ; Skip else block
    
else_block:
    MOVI R1, 0          ; R1 = 0
    
end_if:
    ; Continue...
\end{lstlisting}

\subsection{While Loop}

\begin{lstlisting}[style=assembly, caption={While Loop}]
; int sum = 0;
; int i = 0;
; while (i < 10) {
;     sum += i;
;     i++;
; }

    MOVI R10, 0         ; sum = 0
    MOVI R11, 0         ; i = 0
    
while_start:
    CMPI R11, 10        ; Compare i with 10
    JGE while_end       ; If i >= 10, exit loop
    
    ADD R10, R10, R11   ; sum += i
    ADDI R11, R11, 1    ; i++
    
    JMP while_start     ; Repeat loop
    
while_end:
    ; R10 contains sum
\end{lstlisting}

\subsection{For Loop}

\begin{lstlisting}[style=assembly, caption={For Loop}]
; for (int i = 0; i < 10; i++) {
;     array[i] = i * 2;
; }

    MOVI R10, 0             ; i = 0
    MOVI R11, array         ; R11 = &array[0]
    
for_start:
    CMPI R10, 10            ; i < 10?
    JGE for_end             ; Exit if i >= 10
    
    ; Calculate i * 2
    MULI R12, R10, 2        ; R12 = i * 2
    
    ; Calculate offset
    MULI R13, R10, 8        ; R13 = i * 8 (byte offset)
    ADD R14, R11, R13       ; R14 = &array[i]
    
    ; Store value
    STORE R12, 0(R14)       ; array[i] = i * 2
    
    ADDI R10, R10, 1        ; i++
    JMP for_start
    
for_end:
\end{lstlisting}

% ============================================================================
% CHAPTER 12: SYSTEM CALLS
% ============================================================================
\chapter{System Calls}

\section{System Call Interface}

System calls provide access to operating system services. Arguments are passed in registers R1-R4.

\subsection{System Call Convention}

\begin{itemize}
    \item \textbf{R1:} System call number
    \item \textbf{R2-R4:} Arguments (up to 3)
    \item \textbf{R1:} Return value
\end{itemize}

\begin{lstlisting}[style=assembly]
SYSCALL 0           ; Invoke system call
\end{lstlisting}

\section{Standard System Calls}

Table \ref{tab:syscalls} lists available system calls:

\begin{table}[H]
\centering
\caption{A1M System Calls}
\label{tab:syscalls}
\small
\begin{tabularx}{\textwidth}{|l|l|X|}
\hline
\textbf{Number} & \textbf{Name} & \textbf{Description} \\
\hline
0 & exit & Exit program (R2 = exit code) \\
1 & write & Write to file (R2=fd, R3=buf, R4=len) \\
2 & read & Read from file (R2=fd, R3=buf, R4=len) \\
3 & open & Open file (R2=path, R3=flags) \\
4 & close & Close file (R2=fd) \\
5 & malloc & Allocate memory (R2=size) → R1=address \\
6 & free & Free memory (R2=address) \\
10 & print\_int & Print integer (R2=value) \\
11 & print\_str & Print string (R2=address) \\
12 & read\_int & Read integer → R1=value \\
13 & read\_str & Read string (R2=buffer, R3=max\_len) \\
\hline
\end{tabularx}
\end{table}

\section{System Call Examples}

\subsection{Print String}

\begin{lstlisting}[style=assembly, caption={Print String Example}]
.section .data
message: .string "Hello, World!\n"

.section .text
    MOVI R1, 11         ; syscall: print_str
    MOVI R2, message    ; address of string
    SYSCALL 0           ; invoke
\end{lstlisting}

\subsection{Read Integer}

\begin{lstlisting}[style=assembly, caption={Read Integer Example}]
    MOVI R1, 12         ; syscall: read_int
    SYSCALL 0           ; invoke
    ; R1 now contains the input value
    MOV R10, R1         ; Save to R10
\end{lstlisting}

\subsection{File I/O}

\begin{lstlisting}[style=assembly, caption={File I/O Example}]
.section .data
filename: .string "output.txt"
buffer:   .string "Hello, File!\n"

.section .text
    ; Open file
    MOVI R1, 3          ; syscall: open
    MOVI R2, filename   ; file path
    MOVI R3, 1          ; flags: O_WRONLY
    SYSCALL 0
    MOV R10, R1         ; Save file descriptor
    
    ; Write to file
    MOVI R1, 1          ; syscall: write
    MOV R2, R10         ; file descriptor
    MOVI R3, buffer     ; data to write
    MOVI R4, 13         ; length
    SYSCALL 0
    
    ; Close file
    MOVI R1, 4          ; syscall: close
    MOV R2, R10         ; file descriptor
    SYSCALL 0
\end{lstlisting}

% ============================================================================
% CHAPTER 13: ASSEMBLER IMPLEMENTATION
% ============================================================================
\chapter{Assembler Implementation}

\section{Two-Pass Assembly}

The A1M assembler uses a two-pass algorithm:

\subsection{Pass 1: Symbol Table Construction}

\begin{enumerate}
    \item Read each line of assembly code
    \item Track current address counter
    \item Record labels and their addresses in symbol table
    \item Calculate sizes of data declarations
\end{enumerate}

\subsection{Pass 2: Code Generation}

\begin{enumerate}
    \item Re-read assembly code
    \item Resolve label references using symbol table
    \item Generate machine code for each instruction
    \item Output binary file
\end{enumerate}

\section{Assembler Architecture}

\begin{lstlisting}[style=c, caption={Assembler Main Structure}]
#include <stdio.h>
#include <stdlib.h>
#include <string.h>

#define MAX_LABELS 1000
#define MAX_LINE 256

typedef struct {
    char name[64];
    uint64_t address;
} Symbol;

Symbol symbol_table[MAX_LABELS];
int symbol_count = 0;

// Function prototypes
void pass1(FILE *input);
void pass2(FILE *input, FILE *output);
uint32_t encode_instruction(char *mnemonic, char *operands);
uint64_t lookup_symbol(char *name);

int main(int argc, char **argv) {
    if (argc != 3) {
        fprintf(stderr, "Usage: %s input.asm output.bin\n", 
                argv[0]);
        return 1;
    }
    
    FILE *input = fopen(argv[1], "r");
    if (!input) {
        perror("Error opening input file");
        return 1;
    }
    
    // Pass 1: Build symbol table
    pass1(input);
    rewind(input);
    
    // Pass 2: Generate code
    FILE *output = fopen(argv[2], "wb");
    if (!output) {
        perror("Error opening output file");
        fclose(input);
        return 1;
    }
    
    pass2(input, output);
    
    fclose(input);
    fclose(output);
    
    printf("Assembly successful: %s -> %s\n", 
           argv[1], argv[2]);
    return 0;
}
\end{lstlisting}

\section{Instruction Encoding}

\begin{lstlisting}[style=c, caption={Instruction Encoder}]
uint32_t encode_r_type(uint8_t opcode, uint8_t rd, 
                       uint8_t rs1, uint8_t rs2) {
    uint32_t instruction = 0;
    instruction |= (opcode & 0xFF) << 24;
    instruction |= (rd & 0x1F) << 19;
    instruction |= (rs1 & 0x1F) << 14;
    instruction |= (rs2 & 0x1F) << 9;
    return instruction;
}

uint32_t encode_i_type(uint8_t opcode, uint8_t rd, 
                       uint8_t rs1, int16_t imm) {
    uint32_t instruction = 0;
    instruction |= (opcode & 0xFF) << 24;
    instruction |= (rd & 0x1F) << 19;
    instruction |= (rs1 & 0x1F) << 14;
    instruction |= (imm & 0x3FFF);
    return instruction;
}

uint32_t encode_j_type(uint8_t opcode, uint32_t address) {
    uint32_t instruction = 0;
    instruction |= (opcode & 0xFF) << 24;
    instruction |= (address & 0xFFFFFF);
    return instruction;
}
\end{lstlisting}

\section{Example: Complete Assembler Usage}

\begin{lstlisting}[style=assembly, caption={Complete Assembly Program}]
; Factorial Calculator
; Calculates n! recursively

.section .data
prompt:  .string "Enter number: "
result:  .string "Factorial: "

.section .text
.global _start

_start:
    ; Print prompt
    MOVI R1, 11
    MOVI R2, prompt
    SYSCALL 0
    
    ; Read number
    MOVI R1, 12
    SYSCALL 0
    MOV R4, R1          ; Save input in R4
    
    ; Calculate factorial
    CALL factorial
    
    ; Print result message
    MOVI R1, 11
    MOVI R2, result
    SYSCALL 0
    
    ; Print result value
    MOVI R1, 10
    MOV R2, R1
    SYSCALL 0
    
    ; Exit
    MOVI R1, 0
    MOVI R2, 0
    SYSCALL 0

; Factorial function
; Input: R4 = n
; Output: R1 = n!
factorial:
    PUSH R31            ; Save return address
    PUSH R4             ; Save n
    
    ; Base case: if (n <= 1) return 1
    CMPI R4, 1
    JGT recursive_case
    
    MOVI R1, 1          ; Return 1
    JMP factorial_end
    
recursive_case:
    ; Recursive case: n * factorial(n-1)
    SUBI R4, R4, 1      ; n - 1
    CALL factorial      ; R1 = factorial(n-1)
    
    POP R4              ; Restore original n
    MUL R1, R1, R4      ; R1 = n * factorial(n-1)
    
factorial_end:
    POP R31             ; Restore return address
    RET
\end{lstlisting}

Assemble with:
\begin{verbatim}
$ ./assembler factorial.asm factorial.bin
\end{verbatim}

% ============================================================================
% PART III: EMULATOR
% ============================================================================
\part{Emulator Implementation}

% ============================================================================
% CHAPTER 14: EMULATOR ARCHITECTURE
% ============================================================================
\chapter{Emulator Architecture}

\section{Emulator Overview}

The A1M emulator is a software implementation of the A1M CPU that executes binary programs. It provides a platform for testing and debugging programs before FPGA synthesis.

\subsection{Features}

\begin{itemize}
    \item Cycle-accurate execution
    \item Full instruction set support
    \item System call emulation
    \item Debugging support (breakpoints, single-step)
    \item Memory inspection
    \item Register dump
\end{itemize}

\section{Emulator Core Structure}

\begin{lstlisting}[style=c, caption={Emulator Core}]
#include <stdio.h>
#include <stdint.h>
#include <stdbool.h>
#include <string.h>

#define MEMORY_SIZE (1024 * 1024)  // 1 MB
#define NUM_REGISTERS 32

typedef struct {
    uint64_t registers[NUM_REGISTERS];
    uint64_t pc;                // Program counter
    uint8_t *memory;
    bool zero_flag;
    bool negative_flag;
    bool carry_flag;
    bool overflow_flag;
    bool halted;
    uint64_t cycle_count;
} CPU;

// Initialize CPU
void cpu_init(CPU *cpu) {
    memset(cpu->registers, 0, sizeof(cpu->registers));
    cpu->pc = 0;
    cpu->memory = calloc(MEMORY_SIZE, 1);
    cpu->zero_flag = false;
    cpu->negative_flag = false;
    cpu->carry_flag = false;
    cpu->overflow_flag = false;
    cpu->halted = false;
    cpu->cycle_count = 0;
}

// Load program into memory
bool load_program(CPU *cpu, const char *filename) {
    FILE *f = fopen(filename, "rb");
    if (!f) {
        perror("Error opening file");
        return false;
    }
    
    size_t bytes_read = fread(cpu->memory, 1, 
                              MEMORY_SIZE, f);
    fclose(f);
    
    printf("Loaded %zu bytes\n", bytes_read);
    return true;
}

// Fetch instruction
uint32_t fetch(CPU *cpu) {
    uint32_t instruction = 0;
    memcpy(&instruction, &cpu->memory[cpu->pc], 4);
    cpu->pc += 4;
    return instruction;
}

// Decode instruction
void decode(uint32_t instruction, uint8_t *opcode,
            uint8_t *rd, uint8_t *rs1, uint8_t *rs2,
            int16_t *imm, uint32_t *addr) {
    *opcode = (instruction >> 24) & 0xFF;
    *rd = (instruction >> 19) & 0x1F;
    *rs1 = (instruction >> 14) & 0x1F;
    *rs2 = (instruction >> 9) & 0x1F;
    *imm = instruction & 0x3FFF;
    // Sign extend immediate
    if (*imm & 0x2000) {
        *imm |= 0xC000;
    }
    *addr = instruction & 0xFFFFFF;
}

// Execute one instruction
void execute(CPU *cpu) {
    if (cpu->halted) return;
    
    uint32_t instruction = fetch(cpu);
    
    uint8_t opcode, rd, rs1, rs2;
    int16_t imm;
    uint32_t addr;
    
    decode(instruction, &opcode, &rd, &rs1, &rs2, 
           &imm, &addr);
    
    // Execute based on opcode
    execute_instruction(cpu, opcode, rd, rs1, rs2, 
                       imm, addr);
    
    // R0 always zero
    cpu->registers[0] = 0;
    
    cpu->cycle_count++;
}
\end{lstlisting}

% ============================================================================
% CHAPTER 15: SYSTEM CALL IMPLEMENTATION
% ============================================================================
\chapter{System Call Implementation}

\section{System Call Handler}

The emulator implements system calls to provide I/O and OS functionality during program execution.

\begin{lstlisting}[style=c, caption={System Call Handler}]
void handle_syscall(CPU *cpu) {
    uint64_t syscall_num = cpu->registers[1];
    uint64_t arg1 = cpu->registers[2];
    uint64_t arg2 = cpu->registers[3];
    uint64_t arg3 = cpu->registers[4];
    
    switch (syscall_num) {
        case 0: // exit
            printf("\nProgram exited with code %llu\n", arg1);
            cpu->halted = true;
            break;
            
        case 1: // write
            syscall_write(cpu, arg1, arg2, arg3);
            break;
            
        case 2: // read
            syscall_read(cpu, arg1, arg2, arg3);
            break;
            
        case 10: // print_int
            printf("%lld", (int64_t)arg1);
            fflush(stdout);
            break;
            
        case 11: // print_str
            syscall_print_str(cpu, arg1);
            break;
            
        case 12: // read_int
            {
                int64_t value;
                scanf("%lld", &value);
                cpu->registers[1] = value;
            }
            break;
            
        default:
            fprintf(stderr, "Unknown syscall: %llu\n", 
                   syscall_num);
    }
}
\end{lstlisting}

% ============================================================================
% PART IV: ALGO HIGH-LEVEL LANGUAGE
% ============================================================================
\part{Algo Programming Language}

% ============================================================================
% CHAPTER 16: INTRODUCTION TO ALGO
% ============================================================================
\chapter{Introduction to Algo Language}

\section{Language Overview}

Algo is a C-like, statically-typed programming language designed for the A1M architecture. It implements a comprehensive subset of C features including modern control flow, operators, and type declarations.

\subsection{Design Goals}

\begin{enumerate}[itemsep=8pt]
    \item \textbf{C Compatibility:} Syntax closely follows C language standards
    \item \textbf{Type Safety:} Static type checking with compile-time constants
    \item \textbf{Efficiency:} Direct compilation to A1M assembly (no runtime overhead)
    \item \textbf{Modern Features:} Includes goto, switch/case, do-while, ternary, ++/--
    \item \textbf{Educational:} Clear mapping between high-level and assembly code
\end{enumerate}

\subsection{Hello World}

\begin{lstlisting}[style=algo, caption={Hello World Program}]
// First Algo program
func main() -> void {
    print("Hello, World!\n");
    print("Welcome to Algo programming!\n");
}
\end{lstlisting}

Compilation and execution:
\begin{lstlisting}[style=bash]
$ ./compiler/algo-compiler hello.algo -o hello.asm
$ ./assembler/a1m-asm hello.asm -o hello.bin
$ ./emulator/a1m-emu hello.bin
Hello, World!
Welcome to Algo programming!
\end{lstlisting}

\section{Lexical Structure}

\subsection{Comments}

\begin{lstlisting}[style=algo]
// Single-line comment (C++ style)

/* Multi-line comment
   Can span multiple lines
   C-style block comment */
\end{lstlisting}

\subsection{Identifiers}

Identifiers must start with a letter or underscore, followed by letters, digits, or underscores:

\begin{lstlisting}[style=algo]
let counter: int = 0;        // Valid
let _temp: int = 5;          // Valid
let value123: int = 999;     // Valid
let my_variable: int = 42;   // Valid
// let 123abc: int = 0;      // Invalid - starts with digit
\end{lstlisting}

\subsection{Keywords}

Reserved keywords in Algo (36 C keywords recognized):

\begin{table}[H]
\centering
\caption{Algo Reserved Keywords}
\small
\begin{tabularx}{\textwidth}{|X|X|X|X|}
\hline
\textbf{Types} & \textbf{Control Flow} & \textbf{Operators} & \textbf{Other} \\
\hline
int, long, char & if, else & sizeof & func, return \\
bool, void & while, do, for & ++, -- & let (variable) \\
struct, union & switch, case & +=, -=, *=, /= & enum, typedef \\
signed, unsigned & default, break & \&=, |=, \^{}=, <<= & const, static \\
short, float & continue, goto & >>=, \%=, ?: & extern, register \\
double & & & auto, volatile \\
\hline
\end{tabularx}
\end{table}

\section{Data Types}

\subsection{Primitive Types}

\begin{table}[H]
\centering
\caption{Algo Data Types (64-bit Architecture)}
\begin{tabularx}{\textwidth}{|l|l|l|X|}
\hline
\textbf{Type} & \textbf{Size} & \textbf{Range} & \textbf{Description} \\
\hline
int & 64 bits & $-2^{63}$ to $2^{63}-1$ & Signed integer (default) \\
long & 64 bits & $-2^{63}$ to $2^{63}-1$ & Alias for int \\
char & 8 bits & 0 to 255 & ASCII character \\
bool & 8 bits & 0 or 1 & Boolean (true/false) \\
void & N/A & N/A & Function return type \\
\hline
\end{tabularx}
\end{table}

\textbf{Note:} All types are 64-bit in registers. \code{sizeof(int) = 8}, \code{sizeof(char) = 1}, \code{sizeof(bool) = 1}.

\subsection{Variable Declarations}

Algo uses the \code{let} keyword for variable declarations:

\begin{lstlisting}[style=algo, caption={Variable Declarations}]
// Basic declarations with let keyword
let x: int;                    // Uninitialized
let age: int = 25;             // Initialized
let grade: char = 'A';         // Character
let is_valid: bool = 1;        // Boolean (1 = true, 0 = false)

// Type inference (type required in current implementation)
let count: int = 0;
let message: int = 42;
\end{lstlisting}

\subsection{Constants and Enumerations}

\begin{lstlisting}[style=algo, caption={Compile-Time Constants with Enums}]
// Enum declarations (compile-time constants)
enum Status {
    IDLE,        // Auto: 0
    RUNNING,     // Auto: 1
    PAUSED,      // Auto: 2
    STOPPED      // Auto: 3
}

enum ErrorCode {
    SUCCESS = 0,
    ERROR = -1,
    TIMEOUT = -2
}

// Using enum constants
let status: int = RUNNING;        // status = 1
let err: int = SUCCESS;           // err = 0

// Enum constants are immutable
// RUNNING = 5;  // ERROR: Cannot assign to constant
\end{lstlisting}

\subsection{Literals}

\begin{lstlisting}[style=algo, caption={Literal Values}]
// Integer literals  
let decimal: int = 42;
let negative: int = -100;

// Character literals
let letter: char = 'A';
let newline: char = '\n';
let tab: char = '\t';

// String literals
print("Hello, World\n");

// Boolean values (0 or 1)
let flag1: int = 1;  // true
let flag2: int = 0;  // false
bool flag2 = false;
\end{lstlisting}

\section{Type Conversion}

\subsection{Implicit Conversion}

\begin{lstlisting}[style=algo]
char c = 65;
int x = c;      // char to int (implicit)
// x == 65

bool b = 1;
int y = b;      // bool to int
// y == 1
\end{lstlisting}

\subsection{Explicit Casting}

\begin{lstlisting}[style=algo]
int x = 1000;
char c = (char)x;  // Truncates to 232 (1000 % 256)

int y = 42;
bool b = (bool)y;  // Non-zero becomes true
\end{lstlisting}

\section{Program Structure}

\subsection{Minimal Program}

Every Algo program must have a \code{main} function:

\begin{lstlisting}[style=algo]
fun main() {
    // Program starts here
}
\end{lstlisting}

\subsection{Complete Program with Functions}

\begin{lstlisting}[style=algo, caption={Complete Algo Program Structure}]
// Global variables
int counter = 0;

// Function declaration
fun increment() {
    counter = counter + 1;
}

fun get_counter() -> int {
    return counter;
}

// Main entry point
fun main() {
    increment();
    increment();
    increment();
    
    int value = get_counter();
    print_int(value);  // Prints: 3
}
\end{lstlisting}

% ============================================================================
% CHAPTER 17: EXPRESSIONS AND OPERATORS
% ============================================================================
\chapter{Expressions and Operators}

\section{Arithmetic Operators}

\subsection{Basic Arithmetic}

\begin{table}[H]
\centering
\caption{Arithmetic Operators}
\begin{tabularx}{\textwidth}{|l|l|X|l|}
\hline
\textbf{Operator} & \textbf{Name} & \textbf{Description} & \textbf{Example} \\
\hline
+ & Addition & Add two values & a + b \\
- & Subtraction & Subtract second from first & a - b \\
* & Multiplication & Multiply two values & a * b \\
/ & Division & Integer division & a / b \\
\% & Modulo & Remainder after division & a \% b \\
\hline
\end{tabularx}
\end{table}

\begin{lstlisting}[style=algo, caption={Arithmetic Examples}]
fun main() {
    int a = 10;
    int b = 3;
    
    int sum = a + b;        // 13
    int diff = a - b;       // 7
    int product = a * b;    // 30
    int quotient = a / b;   // 3 (integer division)
    int remainder = a % b;  // 1
    
    // Complex expressions
    int result = (a + b) * 2 - a / b;
    // result = (10 + 3) * 2 - 10 / 3
    // result = 13 * 2 - 3
    // result = 26 - 3 = 23
}
\end{lstlisting}

\subsection{Operator Precedence}

\begin{table}[H]
\centering
\caption{Operator Precedence (Highest to Lowest)}
\begin{tabular}{|l|l|l|}
\hline
\textbf{Precedence} & \textbf{Operators} & \textbf{Associativity} \\
\hline
1 (Highest) & () [] & Left to right \\
2 & ! - (unary) & Right to left \\
3 & * / \% & Left to right \\
4 & + - & Left to right \\
5 & << >> & Left to right \\
6 & < <= > >= & Left to right \\
7 & == != & Left to right \\
8 & \& & Left to right \\
9 & \^{} & Left to right \\
10 & | & Left to right \\
11 & \&\& & Left to right \\
12 (Lowest) & || & Left to right \\
\hline
\end{tabular}
\end{table}

\section{Comparison Operators}

\begin{table}[H]
\centering
\caption{Comparison Operators}
\begin{tabularx}{\textwidth}{|l|l|X|}
\hline
\textbf{Operator} & \textbf{Name} & \textbf{Returns true if} \\
\hline
== & Equal & Operands are equal \\
!= & Not equal & Operands are not equal \\
< & Less than & Left < Right \\
> & Greater than & Left > Right \\
<= & Less or equal & Left $\leq$ Right \\
>= & Greater or equal & Left $\geq$ Right \\
\hline
\end{tabularx}
\end{table}

\begin{lstlisting}[style=algo, caption={Comparison Examples}]
fun main() {
    int x = 10;
    int y = 20;
    
    bool eq = (x == y);     // false
    bool ne = (x != y);     // true
    bool lt = (x < y);      // true
    bool gt = (x > y);      // false
    bool le = (x <= y);     // true
    bool ge = (x >= y);     // false
    
    // Comparisons in conditions
    if (x < y) {
        print("x is less than y\n");
    }
}
\end{lstlisting}

\section{Logical Operators}

\begin{table}[H]
\centering
\caption{Logical Operators}
\begin{tabularx}{\textwidth}{|l|l|X|}
\hline
\textbf{Operator} & \textbf{Name} & \textbf{Description} \\
\hline
\&\& & Logical AND & True if both operands are true \\
|| & Logical OR & True if at least one operand is true \\
! & Logical NOT & Inverts the boolean value \\
\hline
\end{tabularx}
\end{table}

\begin{lstlisting}[style=algo, caption={Logical Operations}]
fun main() {
    bool a = true;
    bool b = false;
    
    bool and_result = a && b;   // false
    bool or_result = a || b;    // true
    bool not_a = !a;            // false
    bool not_b = !b;            // true
    
    // Complex logical expressions
    int x = 10;
    int y = 20;
    int z = 15;
    
    // Check if z is between x and y
    bool in_range = (z > x) && (z < y);  // true
    
    // Check if any condition is true
    bool any_true = (x > y) || (z > x);  // true
}
\end{lstlisting}

\subsection{Short-Circuit Evaluation}

Logical operators use short-circuit evaluation:

\begin{lstlisting}[style=algo]
// AND: If first is false, second is not evaluated
if (x != 0 && y / x > 2) {
    // Safe: division only happens if x != 0
}

// OR: If first is true, second is not evaluated
if (found || search()) {
    // search() not called if found == true
}
\end{lstlisting}

\section{Bitwise Operators}

\begin{table}[H]
\centering
\caption{Bitwise Operators}
\begin{tabularx}{\textwidth}{|l|l|X|}
\hline
\textbf{Operator} & \textbf{Name} & \textbf{Description} \\
\hline
\& & Bitwise AND & Performs AND on each bit \\
| & Bitwise OR & Performs OR on each bit \\
\^{} & Bitwise XOR & Performs XOR on each bit \\
<< & Left shift & Shifts bits left \\
>> & Right shift & Shifts bits right (logical) \\
\hline
\end{tabularx}
\end{table}

\begin{lstlisting}[style=algo, caption={Bitwise Operations}]
fun main() {
    int x = 0b1010;  // 10 in binary
    int y = 0b1100;  // 12 in binary
    
    int and_result = x & y;     // 0b1000 = 8
    int or_result = x | y;      // 0b1110 = 14
    int xor_result = x ^ y;     // 0b0110 = 6
    
    // Shifts
    int left = x << 1;          // 0b10100 = 20
    int right = x >> 1;         // 0b0101 = 5
    
    // Practical uses
    // Set bit 3
    int flags = 0;
    flags = flags | (1 << 3);   // Set bit 3
    
    // Clear bit 3
    flags = flags & ~(1 << 3);  // Clear bit 3
    
    // Toggle bit 3
    flags = flags ^ (1 << 3);   // Toggle bit 3
    
    // Check if bit 3 is set
    bool is_set = (flags & (1 << 3)) != 0;
}
\end{lstlisting}

\section{Assignment Operators}

\begin{lstlisting}[style=algo, caption={Assignment}]
fun main() {
    int x = 10;
    
    // Basic assignment
    x = 20;
    
    // Assignment is an expression (returns the value)
    int y = (x = 30);  // x = 30, y = 30
    
    // Chained assignment
    int a, b, c;
    a = b = c = 5;  // All become 5
}
\end{lstlisting}

\section{Expression Examples}

\subsection{Mathematical Formulas}

\begin{lstlisting}[style=algo, caption={Complex Expressions}]
// Quadratic formula: x = (-b ± sqrt(b²-4ac)) / 2a
fun quadratic_positive(int a, int b, int c) -> int {
    int discriminant = b * b - 4 * a * c;
    // Simplified: no sqrt function yet
    int result = (-b + discriminant) / (2 * a);
    return result;
}

// Area of circle: A = πr²
fun circle_area(int radius) -> int {
    int pi_approx = 314;  // π ≈ 3.14 * 100
    int area = (pi_approx * radius * radius) / 100;
    return area;
}

// Factorial: n! = n × (n-1) × ... × 1
fun factorial(int n) -> int {
    if (n <= 1) {
        return 1;
    }
    return n * factorial(n - 1);
}
\end{lstlisting}

% ============================================================================
% CHAPTER 18: CONTROL STRUCTURES
% ============================================================================
\chapter{Control Structures}

\section{If Statement}

\subsection{Basic If}

\begin{lstlisting}[style=algo, caption={Simple If Statement}]
fun main() {
    int temperature = 25;
    
    if (temperature > 30) {
        print("It's hot!\n");
    }
    
    // Execution continues here
}
\end{lstlisting}

\subsection{If-Else}

\begin{lstlisting}[style=algo, caption={If-Else Statement}]
fun check_even(int n) {
    if (n % 2 == 0) {
        print("Even number\n");
    } else {
        print("Odd number\n");
    }
}

fun main() {
    check_even(42);  // Prints: Even number
    check_even(17);  // Prints: Odd number
}
\end{lstlisting}

\subsection{If-Else-If Chains}

\begin{lstlisting}[style=algo, caption={Multiple Conditions}]
fun grade_to_letter(int score) {
    if (score >= 90) {
        print("Grade: A\n");
    } else if (score >= 80) {
        print("Grade: B\n");
    } else if (score >= 70) {
        print("Grade: C\n");
    } else if (score >= 60) {
        print("Grade: D\n");
    } else {
        print("Grade: F\n");
    }
}

fun main() {
    grade_to_letter(85);  // Prints: Grade: B
    grade_to_letter(55);  // Prints: Grade: F
}
\end{lstlisting}

\subsection{Nested If Statements}

\begin{lstlisting}[style=algo, caption={Nested Conditions}]
fun check_number(int x) {
    if (x > 0) {
        if (x % 2 == 0) {
            print("Positive even\n");
        } else {
            print("Positive odd\n");
        }
    } else if (x < 0) {
        if (x % 2 == 0) {
            print("Negative even\n");
        } else {
            print("Negative odd\n");
        }
    } else {
        print("Zero\n");
    }
}
\end{lstlisting}

\section{While Loop}

\subsection{Basic While Loop}

\begin{lstlisting}[style=algo, caption={While Loop Example}]
fun count_to_ten() {
    int i = 1;
    
    while (i <= 10) {
        print_int(i);
        print(" ");
        i = i + 1;
    }
    print("\n");
}
// Output: 1 2 3 4 5 6 7 8 9 10
\end{lstlisting}

\subsection{While with Complex Condition}

\begin{lstlisting}[style=algo, caption={Sum Until Threshold}]
fun sum_until(int threshold) -> int {
    int sum = 0;
    int i = 1;
    
    while (sum < threshold) {
        sum = sum + i;
        i = i + 1;
    }
    
    return sum;
}

fun main() {
    int result = sum_until(100);
    print_int(result);  // First sum >= 100
}
\end{lstlisting}

\subsection{Infinite Loop with Break (Future Enhancement)}

\begin{lstlisting}[style=algo]
// Note: break/continue not yet implemented
// Current workaround: use conditional flag

fun infinite_loop_example() {
    bool running = true;
    int count = 0;
    
    while (running) {
        count = count + 1;
        if (count >= 10) {
            running = false;  // Exit loop
        }
    }
}
\end{lstlisting}

\section{For Loop}

\subsection{Basic For Loop}

\begin{lstlisting}[style=algo, caption={For Loop Syntax}]
fun print_numbers() {
    // for (initialization; condition; update)
    for (int i = 0; i < 10; i = i + 1) {
        print_int(i);
        print(" ");
    }
    print("\n");
}
// Output: 0 1 2 3 4 5 6 7 8 9
\end{lstlisting}

\subsection{For Loop Variations}

\begin{lstlisting}[style=algo, caption={For Loop Examples}]
// Count down
fun countdown() {
    for (int i = 10; i > 0; i = i - 1) {
        print_int(i);
        print(" ");
    }
    print("Liftoff!\n");
}

// Even numbers
fun print_evens() {
    for (int i = 0; i <= 20; i = i + 2) {
        print_int(i);
        print(" ");
    }
    print("\n");
}

// Powers of 2
fun powers_of_two() {
    for (int i = 1; i <= 1024; i = i * 2) {
        print_int(i);
        print(" ");
    }
    print("\n");
}
\end{lstlisting}

\subsection{Nested Loops}

\begin{lstlisting}[style=algo, caption={Multiplication Table}]
fun multiplication_table() {
    for (int i = 1; i <= 10; i = i + 1) {
        for (int j = 1; j <= 10; j = j + 1) {
            int product = i * j;
            print_int(product);
            print("\t");
        }
        print("\n");
    }
}
\end{lstlisting}

\section{Control Flow Examples}

\subsection{Finding Prime Numbers}

\begin{lstlisting}[style=algo, caption={Prime Number Checker}]
fun is_prime(int n) -> bool {
    if (n <= 1) {
        return false;
    }
    
    if (n == 2) {
        return true;
    }
    
    if (n % 2 == 0) {
        return false;
    }
    
    // Check odd divisors up to sqrt(n)
    int i = 3;
    while (i * i <= n) {
        if (n % i == 0) {
            return false;
        }
        i = i + 2;
    }
    
    return true;
}

fun print_primes(int limit) {
    print("Prime numbers up to ");
    print_int(limit);
    print(":\n");
    
    for (int n = 2; n <= limit; n = n + 1) {
        if (is_prime(n)) {
            print_int(n);
            print(" ");
        }
    }
    print("\n");
}

fun main() {
    print_primes(50);
}
\end{lstlisting}

\subsection{GCD Algorithm}

\begin{lstlisting}[style=algo, caption={Greatest Common Divisor}]
// Euclidean algorithm
fun gcd(int a, int b) -> int {
    while (b != 0) {
        int temp = b;
        b = a % b;
        a = temp;
    }
    return a;
}

fun main() {
    int result = gcd(48, 18);
    print("GCD: ");
    print_int(result);  // Prints: 6
    print("\n");
}
\end{lstlisting}
    print("\n");
}
\end{lstlisting}

% ============================================================================
% CHAPTER 19: FUNCTIONS
% ============================================================================
\chapter{Functions and Procedures}

\section{Function Basics}

\subsection{Function Declaration}

\begin{lstlisting}[style=algo, caption={Function Syntax}]
// Function with no parameters, no return value
fun greet() {
    print("Hello!\n");
}

// Function with parameters, no return
fun print_sum(int a, int b) {
    int sum = a + b;
    print_int(sum);
    print("\n");
}

// Function with return value
fun add(int a, int b) -> int {
    return a + b;
}

// Function with multiple parameters and return
fun max(int a, int b, int c) -> int {
    int result = a;
    
    if (b > result) {
        result = b;
    }
    
    if (c > result) {
        result = c;
    }
    
    return result;
}
\end{lstlisting}

\subsection{Function Calls}

\begin{lstlisting}[style=algo, caption={Calling Functions}]
fun main() {
    // Call function with no return
    greet();
    
    // Call function with parameters
    print_sum(10, 20);
    
    // Use return value
    int sum = add(5, 7);
    print_int(sum);
    
    // Use in expression
    int result = add(3, 4) * 2;
    
    // Nested calls
    int maximum = max(add(1, 2), add(3, 4), add(5, 6));
}
\end{lstlisting}

\section{Recursion}

\subsection{Factorial}

\begin{lstlisting}[style=algo, caption={Recursive Factorial}]
fun factorial(int n) -> int {
    // Base case
    if (n <= 1) {
        return 1;
    }
    
    // Recursive case
    return n * factorial(n - 1);
}

fun main() {
    print("Factorials:\n");
    for (int i = 0; i <= 10; i = i + 1) {
        print_int(i);
        print("! = ");
        print_int(factorial(i));
        print("\n");
    }
}
\end{lstlisting}

\subsection{Fibonacci Sequence}

\begin{lstlisting}[style=algo, caption={Fibonacci Numbers}]
fun fibonacci(int n) -> int {
    if (n <= 1) {
        return n;
    }
    return fibonacci(n - 1) + fibonacci(n - 2);
}

// Iterative version (more efficient)
fun fibonacci_iterative(int n) -> int {
    if (n <= 1) {
        return n;
    }
    
    int prev = 0;
    int curr = 1;
    
    for (int i = 2; i <= n; i = i + 1) {
        int next = prev + curr;
        prev = curr;
        curr = next;
    }
    
    return curr;
}

fun main() {
    print("Fibonacci sequence:\n");
    for (int i = 0; i < 15; i = i + 1) {
        print_int(fibonacci(i));
        print(" ");
    }
    print("\n");
}
\end{lstlisting}

\subsection{Tower of Hanoi}

\begin{lstlisting}[style=algo, caption={Tower of Hanoi Solution}]
fun hanoi(int n, int from, int to, int aux) {
    if (n == 1) {
        print("Move disk 1 from ");
        print_int(from);
        print(" to ");
        print_int(to);
        print("\n");
        return;
    }
    
    hanoi(n - 1, from, aux, to);
    
    print("Move disk ");
    print_int(n);
    print(" from ");
    print_int(from);
    print(" to ");
    print_int(to);
    print("\n");
    
    hanoi(n - 1, aux, to, from);
}

fun main() {
    int disks = 3;
    hanoi(disks, 1, 3, 2);
}
\end{lstlisting}

\section{Scope and Lifetime}

\subsection{Local Variables}

\begin{lstlisting}[style=algo, caption={Variable Scope}]
int global = 100;  // Global scope

fun example() {
    int local = 50;  // Local to example()
    
    print_int(global);  // Can access global
    print_int(local);   // Can access local
}

fun main() {
    print_int(global);  // Can access global
    // print_int(local); // ERROR: local not in scope
    
    int x = 10;  // Local to main()
    {
        int y = 20;  // Nested scope
        print_int(x);  // Can access x
        print_int(y);  // Can access y
    }
    // print_int(y);  // ERROR: y out of scope
}
\end{lstlisting}

% ============================================================================
% CHAPTER 20: ARRAYS
% ============================================================================
\chapter{Arrays and Memory Management}

\section{Array Basics}

\subsection{Array Declaration}

\begin{lstlisting}[style=algo, caption={Array Declaration and Initialization}]
fun main() {
    // Declaration
    int numbers[10];     // Array of 10 integers
    char name[20];       // Array of 20 characters
    bool flags[5];       // Array of 5 booleans
    
    // Initialize elements
    numbers[0] = 10;
    numbers[1] = 20;
    numbers[2] = 30;
    
    // Access elements
    int first = numbers[0];   // 10
    int second = numbers[1];  // 20
}
\end{lstlisting}

\subsection{Array Initialization}

\begin{lstlisting}[style=algo, caption={Initializing Arrays}]
fun initialize_array() {
    int arr[5];
    
    // Initialize with loop
    for (int i = 0; i < 5; i = i + 1) {
        arr[i] = i * 10;
    }
    // arr = {0, 10, 20, 30, 40}
    
    // Print array
    for (int i = 0; i < 5; i = i + 1) {
        print_int(arr[i]);
        print(" ");
    }
    print("\n");
}
\end{lstlisting}

\section{Array Operations}

\subsection{Sum of Array Elements}

\begin{lstlisting}[style=algo, caption={Array Sum}]
fun array_sum(int arr[], int size) -> int {
    int sum = 0;
    
    for (int i = 0; i < size; i = i + 1) {
        sum = sum + arr[i];
    }
    
    return sum;
}

fun main() {
    int numbers[5];
    numbers[0] = 10;
    numbers[1] = 20;
    numbers[2] = 30;
    numbers[3] = 40;
    numbers[4] = 50;
    
    int total = array_sum(numbers, 5);
    print("Sum: ");
    print_int(total);  // Prints: 150
    print("\n");
}
\end{lstlisting}

\subsection{Finding Maximum Element}

\begin{lstlisting}[style=algo, caption={Array Maximum}]
fun array_max(int arr[], int size) -> int {
    if (size == 0) {
        return 0;
    }
    
    int max = arr[0];
    
    for (int i = 1; i < size; i = i + 1) {
        if (arr[i] > max) {
            max = arr[i];
        }
    }
    
    return max;
}

fun array_min(int arr[], int size) -> int {
    if (size == 0) {
        return 0;
    }
    
    int min = arr[0];
    
    for (int i = 1; i < size; i = i + 1) {
        if (arr[i] < min) {
            min = arr[i];
        }
    }
    
    return min;
}
\end{lstlisting}

\subsection{Array Reversal}

\begin{lstlisting}[style=algo, caption={Reverse Array In-Place}]
fun reverse_array(int arr[], int size) {
    int start = 0;
    int end = size - 1;
    
    while (start < end) {
        // Swap elements
        int temp = arr[start];
        arr[start] = arr[end];
        arr[end] = temp;
        
        start = start + 1;
        end = end - 1;
    }
}

fun print_array(int arr[], int size) {
    print("[");
    for (int i = 0; i < size; i = i + 1) {
        print_int(arr[i]);
        if (i < size - 1) {
            print(", ");
        }
    }
    print("]\n");
}

fun main() {
    int numbers[5];
    numbers[0] = 1;
    numbers[1] = 2;
    numbers[2] = 3;
    numbers[3] = 4;
    numbers[4] = 5;
    
    print("Original: ");
    print_array(numbers, 5);
    
    reverse_array(numbers, 5);
    
    print("Reversed: ");
    print_array(numbers, 5);
}
\end{lstlisting}

\section{Sorting Algorithms}

\subsection{Bubble Sort}

\begin{lstlisting}[style=algo, caption={Bubble Sort Implementation}]
fun bubble_sort(int arr[], int n) {
    for (int i = 0; i < n - 1; i = i + 1) {
        for (int j = 0; j < n - i - 1; j = j + 1) {
            if (arr[j] > arr[j + 1]) {
                // Swap
                int temp = arr[j];
                arr[j] = arr[j + 1];
                arr[j + 1] = temp;
            }
        }
    }
}

fun main() {
    int data[8];
    data[0] = 64;
    data[1] = 34;
    data[2] = 25;
    data[3] = 12;
    data[4] = 22;
    data[5] = 11;
    data[6] = 90;
    data[7] = 88;
    
    print("Before sorting:\n");
    print_array(data, 8);
    
    bubble_sort(data, 8);
    
    print("After sorting:\n");
    print_array(data, 8);
}
\end{lstlisting}

\subsection{Selection Sort}

\begin{lstlisting}[style=algo, caption={Selection Sort}]
fun selection_sort(int arr[], int n) {
    for (int i = 0; i < n - 1; i = i + 1) {
        int min_idx = i;
        
        // Find minimum element
        for (int j = i + 1; j < n; j = j + 1) {
            if (arr[j] < arr[min_idx]) {
                min_idx = j;
            }
        }
        
        // Swap minimum with first unsorted element
        if (min_idx != i) {
            int temp = arr[i];
            arr[i] = arr[min_idx];
            arr[min_idx] = temp;
        }
    }
}
\end{lstlisting}

\subsection{Insertion Sort}

\begin{lstlisting}[style=algo, caption={Insertion Sort}]
fun insertion_sort(int arr[], int n) {
    for (int i = 1; i < n; i = i + 1) {
        int key = arr[i];
        int j = i - 1;
        
        // Move elements greater than key
        while (j >= 0 && arr[j] > key) {
            arr[j + 1] = arr[j];
            j = j - 1;
        }
        
        arr[j + 1] = key;
    }
}
\end{lstlisting}

\section{Searching Algorithms}

\subsection{Linear Search}

\begin{lstlisting}[style=algo, caption={Linear Search}]
fun linear_search(int arr[], int n, int target) -> int {
    for (int i = 0; i < n; i = i + 1) {
        if (arr[i] == target) {
            return i;  // Found at index i
        }
    }
    return -1;  // Not found
}

fun main() {
    int numbers[5];
    numbers[0] = 10;
    numbers[1] = 20;
    numbers[2] = 30;
    numbers[3] = 40;
    numbers[4] = 50;
    
    int index = linear_search(numbers, 5, 30);
    if (index != -1) {
        print("Found at index: ");
        print_int(index);
    } else {
        print("Not found");
    }
    print("\n");
}
\end{lstlisting}

\subsection{Binary Search}

\begin{lstlisting}[style=algo, caption={Binary Search (Requires Sorted Array)}]
fun binary_search(int arr[], int n, int target) -> int {
    int left = 0;
    int right = n - 1;
    
    while (left <= right) {
        int mid = (left + right) / 2;
        
        if (arr[mid] == target) {
            return mid;  // Found
        }
        
        if (arr[mid] < target) {
            left = mid + 1;  // Search right half
        } else {
            right = mid - 1;  // Search left half
        }
    }
    
    return -1;  // Not found
}
\end{lstlisting}

\section{String Operations}

\subsection{String Length}

\begin{lstlisting}[style=algo, caption={String Functions}]
fun string_length(char str[]) -> int {
    int len = 0;
    
    while (str[len] != '\0') {
        len = len + 1;
    }
    
    return len;
}

fun string_copy(char dest[], char src[]) {
    int i = 0;
    
    while (src[i] != '\0') {
        dest[i] = src[i];
        i = i + 1;
    }
    
    dest[i] = '\0';  // Null terminate
}

fun string_compare(char str1[], char str2[]) -> bool {
    int i = 0;
    
    while (str1[i] != '\0' && str2[i] != '\0') {
        if (str1[i] != str2[i]) {
            return false;
        }
        i = i + 1;
    }
    
    return str1[i] == str2[i];
}
\end{lstlisting}

% ============================================================================
% CHAPTER 21: ADVANCED FEATURES
% ============================================================================
\chapter{Advanced Features}

\section{Multi-Dimensional Arrays}

\subsection{2D Arrays (Future Enhancement)}

\begin{lstlisting}[style=algo]
// Note: True 2D arrays not yet supported
// Workaround: Use 1D array with manual indexing

fun get_2d_index(int row, int col, int cols) -> int {
    return row * cols + col;
}

fun matrix_example() {
    int matrix[9];  // 3x3 matrix
    int rows = 3;
    int cols = 3;
    
    // Initialize
    for (int i = 0; i < rows; i = i + 1) {
        for (int j = 0; j < cols; j = j + 1) {
            int index = get_2d_index(i, j, cols);
            matrix[index] = i * cols + j;
        }
    }
    
    // Access element at [1][2]
    int idx = get_2d_index(1, 2, cols);
    int value = matrix[idx];
}
\end{lstlisting}

\section{Pointer-Like Operations (Future)}

\begin{lstlisting}[style=algo]
// Future enhancement: Pointer support
// Current: Arrays passed by reference

fun modify_array(int arr[], int size) {
    // Arrays are passed by reference
    // Changes affect original array
    for (int i = 0; i < size; i = i + 1) {
        arr[i] = arr[i] * 2;
    }
}
\end{lstlisting}

% ============================================================================
% CHAPTER 22: STANDARD LIBRARY
% ============================================================================
\chapter{Standard Library Functions}

\section{I/O Functions}

\begin{table}[H]
\centering
\caption{Standard I/O Functions}
\begin{tabularx}{\textwidth}{|l|l|X|}
\hline
\textbf{Function} & \textbf{Parameters} & \textbf{Description} \\
\hline
print & string & Print string to stdout \\
print\_int & int & Print integer value \\
print\_char & char & Print single character \\
read\_int & - & Read integer from stdin \\
read\_char & - & Read character from stdin \\
\hline
\end{tabularx}
\end{table}

\begin{lstlisting}[style=algo, caption={Using I/O Functions}]
fun main() {
    print("Enter your age: ");
    int age = read_int();
    
    print("You are ");
    print_int(age);
    print(" years old\n");
    
    print("Enter initial: ");
    char initial = read_char();
    
    print("Your initial is: ");
    print_char(initial);
    print("\n");
}
\end{lstlisting}

% ============================================================================
% CHAPTER 23: COMPLETE ALGO PROGRAMS
% ============================================================================
\chapter{Complete Algo Programs}

\section{Calculator Program}

\begin{lstlisting}[style=algo, caption={Simple Calculator}]
fun calculator() {
    print("Simple Calculator\n");
    print("================\n");
    
    print("Enter first number: ");
    int a = read_int();
    
    print("Enter operation (+,-,*,/): ");
    char op = read_char();
    
    print("Enter second number: ");
    int b = read_int();
    
    int result = 0;
    bool valid = true;
    
    if (op == '+') {
        result = a + b;
    } else if (op == '-') {
        result = a - b;
    } else if (op == '*') {
        result = a * b;
    } else if (op == '/') {
        if (b != 0) {
            result = a / b;
        } else {
            print("Error: Division by zero\n");
            valid = false;
        }
    } else {
        print("Error: Invalid operation\n");
        valid = false;
    }
    
    if (valid) {
        print("Result: ");
        print_int(result);
        print("\n");
    }
}

fun main() {
    calculator();
}
\end{lstlisting}

\section{Number Guessing Game}

\begin{lstlisting}[style=algo, caption={Guess the Number}]
fun guessing_game() {
    int secret = 42;  // In real version, use random number
    int attempts = 0;
    int max_attempts = 10;
    bool won = false;
    
    print("Guess the number (1-100)\n");
    print("You have ");
    print_int(max_attempts);
    print(" attempts\n\n");
    
    while (attempts < max_attempts && !won) {
        print("Attempt ");
        print_int(attempts + 1);
        print(": ");
        
        int guess = read_int();
        attempts = attempts + 1;
        
        if (guess == secret) {
            print("Congratulations! You won!\n");
            won = true;
        } else if (guess < secret) {
            print("Too low!\n");
        } else {
            print("Too high!\n");
        }
    }
    
    if (!won) {
        print("Game over! The number was ");
        print_int(secret);
        print("\n");
    }
}

fun main() {
    guessing_game();
}
\end{lstlisting}

\section{Student Grade Management}

\begin{lstlisting}[style=algo, caption={Grade Management System}]
fun calculate_average(int scores[], int count) -> int {
    int sum = 0;
    
    for (int i = 0; i < count; i = i + 1) {
        sum = sum + scores[i];
    }
    
    return sum / count;
}

fun find_grade(int score) -> char {
    if (score >= 90) {
        return 'A';
    } else if (score >= 80) {
        return 'B';
    } else if (score >= 70) {
        return 'C';
    } else if (score >= 60) {
        return 'D';
    } else {
        return 'F';
    }
}

fun grade_system() {
    int num_students = 5;
    int scores[5];
    
    print("Enter scores for ");
    print_int(num_students);
    print(" students:\n");
    
    for (int i = 0; i < num_students; i = i + 1) {
        print("Student ");
        print_int(i + 1);
        print(": ");
        scores[i] = read_int();
    }
    
    print("\nResults:\n");
    print("========\n");
    
    for (int i = 0; i < num_students; i = i + 1) {
        print("Student ");
        print_int(i + 1);
        print(": ");
        print_int(scores[i]);
        print(" (Grade: ");
        print_char(find_grade(scores[i]));
        print(")\n");
    }
    
    int avg = calculate_average(scores, num_students);
    print("\nClass Average: ");
    print_int(avg);
    print("\n");
}

fun main() {
    grade_system();
}
\end{lstlisting}

\section{Matrix Operations}

\begin{lstlisting}[style=algo, caption={Matrix Addition}]
fun matrix_add(int a[], int b[], int result[], 
               int rows, int cols) {
    int size = rows * cols;
    
    for (int i = 0; i < size; i = i + 1) {
        result[i] = a[i] + b[i];
    }
}

fun matrix_print(int matrix[], int rows, int cols) {
    for (int i = 0; i < rows; i = i + 1) {
        for (int j = 0; j < cols; j = j + 1) {
            int index = i * cols + j;
            print_int(matrix[index]);
            print("\t");
        }
        print("\n");
    }
}

fun main() {
    int rows = 2;
    int cols = 2;
    int size = rows * cols;
    
    int matrix_a[4];
    int matrix_b[4];
    int result[4];
    
    // Initialize matrices
    matrix_a[0] = 1; matrix_a[1] = 2;
    matrix_a[2] = 3; matrix_a[3] = 4;
    
    matrix_b[0] = 5; matrix_b[1] = 6;
    matrix_b[2] = 7; matrix_b[3] = 8;
    
    // Add matrices
    matrix_add(matrix_a, matrix_b, result, rows, cols);
    
    print("Matrix A:\n");
    matrix_print(matrix_a, rows, cols);
    
    print("\nMatrix B:\n");
    matrix_print(matrix_b, rows, cols);
    
    print("\nA + B:\n");
    matrix_print(result, rows, cols);
}
\end{lstlisting}

% ============================================================================
% CHAPTER 24: ARRAYS (Detailed Implementation)
% ============================================================================
\chapter{Arrays in Algo}

\section{Array Declaration and Initialization}

Arrays in Algo are statically-sized, homogeneous collections allocated on the stack or in static memory.

\subsection{Basic Array Declaration}

\begin{lstlisting}[style=algo, caption={Array Declaration Syntax}]
// Fixed-size integer array
int numbers[10];

// Character array (string)
char message[256];

// Boolean array
bool flags[8];

// Long integer array
long bigNumbers[100];
\end{lstlisting}

\subsection{Array Initialization}

Currently, Algo supports declaration with size specification. Manual initialization is required:

\begin{lstlisting}[style=algo, caption={Initializing Arrays}]
fun init_array() {
    int arr[5];
    
    // Manual initialization
    arr[0] = 10;
    arr[1] = 20;
    arr[2] = 30;
    arr[3] = 40;
    arr[4] = 50;
    
    return 0;
}
\end{lstlisting}

\textbf{Note:} Array initializer lists (e.g., \code{int arr[5] = \{1, 2, 3, 4, 5\}}) are \textbf{not yet implemented}.

\section{Array Access and Manipulation}

\subsection{Indexing}

\begin{lstlisting}[style=algo, caption={Array Element Access}]
fun array_operations() {
    int data[10];
    int i = 0;
    
    // Write elements
    while (i < 10) {
        data[i] = i * i;  // Store squares
        i = i + 1;
    }
    
    // Read and print elements
    i = 0;
    while (i < 10) {
        print_int(data[i]);
        print_char('\n');
        i = i + 1;
    }
}
\end{lstlisting}

\subsection{Bounds Checking}

\textbf{Current Implementation:} Algo compiler does \textbf{NOT} perform automatic bounds checking.

\begin{lstlisting}[style=algo, caption={Array Bounds - Manual Checking Required}]
fun safe_access(int arr[], int size, int index) -> int {
    // Manual bounds check
    if (index < 0 || index >= size) {
        print("Error: Index out of bounds\n");
        return -1;
    }
    return arr[index];
}
\end{lstlisting}

\section{Arrays as Function Parameters}

Arrays are passed by reference (pointer semantics):

\begin{lstlisting}[style=algo, caption={Passing Arrays to Functions}]
// Function accepts array parameter
fun sum_array(int arr[], int size) -> int {
    int sum = 0;
    int i = 0;
    
    while (i < size) {
        sum = sum + arr[i];
        i = i + 1;
    }
    
    return sum;
}

fun main() {
    int numbers[5];
    numbers[0] = 10;
    numbers[1] = 20;
    numbers[2] = 30;
    numbers[3] = 40;
    numbers[4] = 50;
    
    int total = sum_array(numbers, 5);
    print_int(total);  // Prints: 150
}
\end{lstlisting}

\section{String Handling}

Strings are character arrays with null termination:

\begin{lstlisting}[style=algo, caption={String Operations}]
fun string_length(char str[]) -> int {
    int len = 0;
    while (str[len] != '\0') {
        len = len + 1;
    }
    return len;
}

fun string_copy(char dest[], char src[]) {
    int i = 0;
    while (src[i] != '\0') {
        dest[i] = src[i];
        i = i + 1;
    }
    dest[i] = '\0';  // Null terminator
}

fun main() {
    char greeting[50];
    char message[50] = "Hello, World!";  // String literal
    
    string_copy(greeting, message);
    print_str(greeting);
}
\end{lstlisting}

\textbf{Note:} String literals are supported, but string concatenation operators are \textbf{not yet implemented}.

\section{Array Implementation Details}

\subsection{Memory Layout}

\begin{verbatim}
Stack Layout for: int arr[5];

Higher addresses
    ┌─────────────┐
    │   arr[4]    │  SP + 32
    ├─────────────┤
    │   arr[3]    │  SP + 24
    ├─────────────┤
    │   arr[2]    │  SP + 16
    ├─────────────┤
    │   arr[1]    │  SP + 8
    ├─────────────┤
    │   arr[0]    │  SP + 0  (base address)
    └─────────────┘
Lower addresses
\end{verbatim}

Each element occupies 8 bytes (64-bit architecture).

\subsection{Assembly Code Generation}

\begin{lstlisting}[style=assembly, caption={Generated Assembly for Array Access}]
; Algo: arr[i] = 42;
; Assume arr base at R10, i in R11

    MULI R12, R11, 8     ; R12 = i * 8 (element size)
    ADD R13, R10, R12    ; R13 = arr + (i * 8)
    MOVI R14, 42         ; R14 = value
    STORE R14, R13, 0    ; *R13 = 42
\end{lstlisting}

% ============================================================================
% CHAPTER 25: ADVANCED FEATURES
% ============================================================================
\chapter{Advanced Features}

\section{Implemented Advanced Features}

\subsection{Nested Function Calls}

\begin{lstlisting}[style=algo, caption={Nested Function Calls}]
fun factorial(int n) -> int {
    if (n <= 1) {
        return 1;
    }
    return n * factorial(n - 1);  // Recursive call
}

fun compute(int a, int b) -> int {
    return factorial(a) + factorial(b);  // Nested calls
}

fun main() {
    int result = compute(5, 3);  // 120 + 6 = 126
    print_int(result);
}
\end{lstlisting}

\subsection{Complex Expressions}

\begin{lstlisting}[style=algo, caption={Multi-Level Expressions}]
fun evaluate(int x, int y, int z) -> int {
    // Operator precedence and associativity
    return (x + y * z) - (x / y) + (z % x);
}

fun logical_ops(bool a, bool b, bool c) -> bool {
    // Short-circuit evaluation
    return (a && b) || (!c && (a || b));
}
\end{lstlisting}

\section{Partially Implemented Features}

\subsection{Multi-Dimensional Arrays}

\textbf{Status:} \textcolor{red}{NOT IMPLEMENTED}

Workaround using single-dimensional arrays:

\begin{lstlisting}[style=algo, caption={Simulating 2D Arrays}]
fun get_2d(int arr[], int rows, int cols, int r, int c) -> int {
    // Manual 2D indexing: arr[r][c] = arr[r * cols + c]
    int index = r * cols + c;
    return arr[index];
}

fun set_2d(int arr[], int rows, int cols, int r, int c, int val) {
    int index = r * cols + c;
    arr[index] = val;
}

fun main() {
    int matrix[9];  // 3x3 matrix as flat array
    
    set_2d(matrix, 3, 3, 0, 0, 1);  // matrix[0][0] = 1
    set_2d(matrix, 3, 3, 1, 1, 5);  // matrix[1][1] = 5
    
    int val = get_2d(matrix, 3, 3, 0, 0);
    print_int(val);
}
\end{lstlisting}

\subsection{Structures and User-Defined Types}

\textbf{Status:} \textcolor{red}{NOT IMPLEMENTED}

Planned syntax (future):

\begin{lstlisting}[style=algo, caption={Structures (Planned Feature)}]
// NOT YET SUPPORTED - Future enhancement
struct Point {
    int x;
    int y;
}

struct Rectangle {
    Point topLeft;
    Point bottomRight;
}

fun area(Rectangle rect) -> int {
    int width = rect.bottomRight.x - rect.topLeft.x;
    int height = rect.topLeft.y - rect.bottomRight.y;
    return width * height;
}
\end{lstlisting}

\subsection{Pointers and Dynamic Memory}

\textbf{Status:} \textcolor{red}{PARTIALLY IMPLEMENTED}

Arrays are passed by reference, but explicit pointer syntax is not supported:

\begin{lstlisting}[style=algo, caption={Current Array Reference Behavior}]
fun modify_array(int arr[]) {
    arr[0] = 100;  // Modifies original array
}

fun main() {
    int data[5];
    data[0] = 10;
    
    modify_array(data);  // Passes reference
    print_int(data[0]);  // Prints: 100
}
\end{lstlisting}

Explicit pointer syntax (e.g., \code{int* ptr}, \code{*ptr}, \code{\&var}) is \textbf{NOT IMPLEMENTED}.

\section{Unimplemented Advanced Features}

\begin{table}[H]
\centering
\caption{Feature Implementation Status}
\begin{tabularx}{\textwidth}{|X|l|X|}
\hline
\textbf{Feature} & \textbf{Status} & \textbf{Notes} \\
\hline
Multi-dimensional arrays & \textcolor{red}{No} & Use flat arrays with manual indexing \\
\hline
Structures/Records & \textcolor{red}{No} & Planned for future versions \\
\hline
Explicit pointers & \textcolor{red}{No} & Arrays use implicit pointers \\
\hline
Dynamic memory (malloc/free) & \textcolor{red}{No} & Only stack allocation \\
\hline
Enumerations & \textcolor{red}{No} & Use integer constants \\
\hline
Type aliases (typedef) & \textcolor{red}{No} & - \\
\hline
Function pointers & \textcolor{red}{No} & - \\
\hline
Variable-length arrays & \textcolor{red}{No} & Only fixed-size arrays \\
\hline
Bit fields & \textcolor{red}{No} & - \\
\hline
\end{tabularx}
\end{table}

% ============================================================================
% CHAPTER 26: STANDARD LIBRARY
% ============================================================================
\chapter{Standard Library}

\section{I/O Functions}

\subsection{Console Output}

\begin{lstlisting}[style=algo, caption={Standard Output Functions}]
// Print integer value
fun print_int(int value);

// Print character
fun print_char(char c);

// Print string (null-terminated char array)
fun print_str(char str[]);

// Example usage
fun main() {
    print_str("The answer is: ");
    print_int(42);
    print_char('\n');
}
\end{lstlisting}

\subsection{Console Input}

\begin{lstlisting}[style=algo, caption={Standard Input Functions}]
// Read integer from stdin
fun read_int() -> int;

// Read single character
fun read_char() -> char;

// Read string into buffer
fun read_str(char buffer[], int maxlen) -> int;

// Example usage
fun main() {
    print_str("Enter a number: ");
    int num = read_int();
    
    print_str("You entered: ");
    print_int(num);
    print_char('\n');
}
\end{lstlisting}

\section{String Utility Functions}

\textbf{Status:} \textcolor{orange}{User must implement - not in standard library}

Common string functions that users need to implement:

\begin{lstlisting}[style=algo, caption={User-Implemented String Functions}]
// String length
fun strlen(char str[]) -> int {
    int len = 0;
    while (str[len] != '\0') {
        len = len + 1;
    }
    return len;
}

// String comparison
fun strcmp(char s1[], char s2[]) -> int {
    int i = 0;
    while (s1[i] != '\0' && s2[i] != '\0') {
        if (s1[i] != s2[i]) {
            return s1[i] - s2[i];
        }
        i = i + 1;
    }
    return s1[i] - s2[i];
}

// String concatenation
fun strcat(char dest[], char src[]) {
    int dest_len = strlen(dest);
    int i = 0;
    while (src[i] != '\0') {
        dest[dest_len + i] = src[i];
        i = i + 1;
    }
    dest[dest_len + i] = '\0';
}
\end{lstlisting}

\section{Mathematical Functions}

\textbf{Status:} \textcolor{orange}{Basic operators only}

Currently supported:
\begin{itemize}
    \item Arithmetic: \code{+}, \code{-}, \code{*}, \code{/}, \code{\%}
    \item Comparison: \code{<}, \code{>}, \code{<=}, \code{>=}, \code{==}, \code{!=}
    \item Logical: \code{\&\&}, \code{||}, \code{!}
\end{itemize}

\textbf{Not implemented:} Transcendental functions (sin, cos, sqrt, pow, etc.)

\begin{lstlisting}[style=algo, caption={User-Implemented Math Functions}]
// Absolute value
fun abs(int x) -> int {
    if (x < 0) {
        return -x;
    }
    return x;
}

// Integer power
fun power(int base, int exp) -> int {
    int result = 1;
    int i = 0;
    while (i < exp) {
        result = result * base;
        i = i + 1;
    }
    return result;
}

// Square root (integer approximation)
fun sqrt(int n) -> int {
    if (n < 0) return -1;
    if (n == 0) return 0;
    
    int x = n;
    int y = (x + 1) / 2;
    
    while (y < x) {
        x = y;
        y = (x + n / x) / 2;
    }
    return x;
}
\end{lstlisting}

\section{System Call Interface}

\subsection{Exit Function}

\begin{lstlisting}[style=algo, caption={Program Termination}]
// Exit with status code
fun exit(int code);

fun main() {
    print_str("Program complete\n");
    exit(0);  // Exit with success
}
\end{lstlisting}

\subsection{Available System Calls}

\begin{table}[H]
\centering
\caption{Implemented System Calls}
\small
\begin{tabularx}{\textwidth}{|l|l|X|}
\hline
\textbf{Function} & \textbf{Signature} & \textbf{Description} \\
\hline
exit & \code{fun exit(int code)} & Terminate program \\
\hline
print\_int & \code{fun print\_int(int val)} & Print integer \\
\hline
print\_char & \code{fun print\_char(char c)} & Print character \\
\hline
print\_str & \code{fun print\_str(char s[])} & Print string \\
\hline
read\_int & \code{fun read\_int() -> int} & Read integer \\
\hline
read\_char & \code{fun read\_char() -> char} & Read character \\
\hline
read\_str & \code{fun read\_str(char b[], int n) -> int} & Read string \\
\hline
\end{tabularx}
\end{table}

\textbf{Not implemented:} File I/O, memory allocation (malloc/free), timing functions, random number generation.

% ============================================================================
% CHAPTER 27: COMPLETE ALGO PROGRAMS
% ============================================================================
\chapter{Complete Algo Programs}

\section{Sorting Algorithms}

\subsection{Bubble Sort}

\begin{lstlisting}[style=algo, caption={Bubble Sort Implementation}]
fun bubble_sort(int arr[], int n) {
    int i = 0;
    int j = 0;
    int temp = 0;
    
    while (i < n - 1) {
        j = 0;
        while (j < n - i - 1) {
            if (arr[j] > arr[j + 1]) {
                // Swap
                temp = arr[j];
                arr[j] = arr[j + 1];
                arr[j + 1] = temp;
            }
            j = j + 1;
        }
        i = i + 1;
    }
}

fun print_array(int arr[], int n) {
    int i = 0;
    while (i < n) {
        print_int(arr[i]);
        print_char(' ');
        i = i + 1;
    }
    print_char('\n');
}

fun main() {
    int numbers[8];
    numbers[0] = 64;
    numbers[1] = 34;
    numbers[2] = 25;
    numbers[3] = 12;
    numbers[4] = 22;
    numbers[5] = 11;
    numbers[6] = 90;
    numbers[7] = 88;
    
    print_str("Before sorting:\n");
    print_array(numbers, 8);
    
    bubble_sort(numbers, 8);
    
    print_str("After sorting:\n");
    print_array(numbers, 8);
}
\end{lstlisting}

\subsection{Binary Search}

\begin{lstlisting}[style=algo, caption={Binary Search Implementation}]
fun binary_search(int arr[], int n, int target) -> int {
    int left = 0;
    int right = n - 1;
    int mid = 0;
    
    while (left <= right) {
        mid = (left + right) / 2;
        
        if (arr[mid] == target) {
            return mid;  // Found
        }
        
        if (arr[mid] < target) {
            left = mid + 1;
        } else {
            right = mid - 1;
        }
    }
    
    return -1;  // Not found
}

fun main() {
    int sorted[10];
    sorted[0] = 1;
    sorted[1] = 3;
    sorted[2] = 5;
    sorted[3] = 7;
    sorted[4] = 9;
    sorted[5] = 11;
    sorted[6] = 13;
    sorted[7] = 15;
    sorted[8] = 17;
    sorted[9] = 19;
    
    int index = binary_search(sorted, 10, 13);
    
    if (index >= 0) {
        print_str("Found at index: ");
        print_int(index);
        print_char('\n');
    } else {
        print_str("Not found\n");
    }
}
\end{lstlisting}

\section{Mathematical Programs}

\subsection{Fibonacci Sequence}

\begin{lstlisting}[style=algo, caption={Fibonacci Generator}]
// Recursive implementation
fun fibonacci_recursive(int n) -> int {
    if (n <= 1) {
        return n;
    }
    return fibonacci_recursive(n - 1) + fibonacci_recursive(n - 2);
}

// Iterative implementation (more efficient)
fun fibonacci_iterative(int n) -> int {
    if (n <= 1) return n;
    
    int prev = 0;
    int curr = 1;
    int next = 0;
    int i = 2;
    
    while (i <= n) {
        next = prev + curr;
        prev = curr;
        curr = next;
        i = i + 1;
    }
    
    return curr;
}

fun main() {
    int n = 10;
    int i = 0;
    
    print_str("First 10 Fibonacci numbers:\n");
    while (i < n) {
        print_int(fibonacci_iterative(i));
        print_char(' ');
        i = i + 1;
    }
    print_char('\n');
}
\end{lstlisting}

\subsection{Prime Number Checker}

\begin{lstlisting}[style=algo, caption={Prime Number Detection}]
fun is_prime(int n) -> bool {
    if (n <= 1) return false;
    if (n <= 3) return true;
    
    if (n % 2 == 0 || n % 3 == 0) {
        return false;
    }
    
    int i = 5;
    while (i * i <= n) {
        if (n % i == 0 || n % (i + 2) == 0) {
            return false;
        }
        i = i + 6;
    }
    
    return true;
}

fun print_primes(int limit) {
    int num = 2;
    int count = 0;
    
    print_str("Prime numbers up to ");
    print_int(limit);
    print_str(":\n");
    
    while (num <= limit) {
        if (is_prime(num)) {
            print_int(num);
            print_char(' ');
            count = count + 1;
            
            if (count % 10 == 0) {
                print_char('\n');
            }
        }
        num = num + 1;
    }
    
    print_char('\n');
    print_str("Total: ");
    print_int(count);
    print_char('\n');
}

fun main() {
    print_primes(100);
}
\end{lstlisting}

\section{Data Structure Examples}

\subsection{Stack Implementation}

\begin{lstlisting}[style=algo, caption={Stack Using Arrays}]
// Global stack variables (Algo doesn't have structures)
int stack[100];
int stack_top = -1;

fun stack_push(int value) -> bool {
    if (stack_top >= 99) {
        print_str("Stack overflow!\n");
        return false;
    }
    
    stack_top = stack_top + 1;
    stack[stack_top] = value;
    return true;
}

fun stack_pop() -> int {
    if (stack_top < 0) {
        print_str("Stack underflow!\n");
        return -1;
    }
    
    int value = stack[stack_top];
    stack_top = stack_top - 1;
    return value;
}

fun stack_peek() -> int {
    if (stack_top < 0) {
        return -1;
    }
    return stack[stack_top];
}

fun stack_is_empty() -> bool {
    return stack_top < 0;
}

fun main() {
    stack_push(10);
    stack_push(20);
    stack_push(30);
    
    print_str("Top element: ");
    print_int(stack_peek());
    print_char('\n');
    
    while (!stack_is_empty()) {
        print_int(stack_pop());
        print_char(' ');
    }
    print_char('\n');
}
\end{lstlisting}

\subsection{Matrix Operations}

\begin{lstlisting}[style=algo, caption={Matrix Multiplication}]
fun matrix_multiply(int a[], int b[], int result[], 
                    int rows_a, int cols_a, int cols_b) {
    int i = 0;
    int j = 0;
    int k = 0;
    int sum = 0;
    
    while (i < rows_a) {
        j = 0;
        while (j < cols_b) {
            sum = 0;
            k = 0;
            while (k < cols_a) {
                sum = sum + a[i * cols_a + k] * b[k * cols_b + j];
                k = k + 1;
            }
            result[i * cols_b + j] = sum;
            j = j + 1;
        }
        i = i + 1;
    }
}

fun print_matrix(int mat[], int rows, int cols) {
    int i = 0;
    int j = 0;
    
    while (i < rows) {
        j = 0;
        while (j < cols) {
            print_int(mat[i * cols + j]);
            print_char(' ');
            j = j + 1;
        }
        print_char('\n');
        i = i + 1;
    }
}

fun main() {
    int a[6];  // 2x3 matrix
    int b[6];  // 3x2 matrix
    int c[4];  // 2x2 result
    
    // Initialize matrix A (2x3)
    a[0] = 1; a[1] = 2; a[2] = 3;
    a[3] = 4; a[4] = 5; a[5] = 6;
    
    // Initialize matrix B (3x2)
    b[0] = 7; b[1] = 8;
    b[2] = 9; b[3] = 10;
    b[4] = 11; b[5] = 12;
    
    matrix_multiply(a, b, c, 2, 3, 2);
    
    print_str("Result (2x2):\n");
    print_matrix(c, 2, 2);
}
\end{lstlisting}

\section{Complete Application Example}

\begin{lstlisting}[style=algo, caption={Simple Calculator Program}]
fun add(int a, int b) -> int { return a + b; }
fun subtract(int a, int b) -> int { return a - b; }
fun multiply(int a, int b) -> int { return a * b; }
fun divide(int a, int b) -> int {
    if (b == 0) {
        print_str("Error: Division by zero\n");
        return 0;
    }
    return a / b;
}

fun print_menu() {
    print_str("\n=== Simple Calculator ===\n");
    print_str("1. Add\n");
    print_str("2. Subtract\n");
    print_str("3. Multiply\n");
    print_str("4. Divide\n");
    print_str("5. Exit\n");
    print_str("Enter choice: ");
}

fun main() {
    int choice = 0;
    int num1 = 0;
    int num2 = 0;
    int result = 0;
    bool running = true;
    
    while (running) {
        print_menu();
        choice = read_int();
        
        if (choice == 5) {
            print_str("Goodbye!\n");
            running = false;
        } else if (choice >= 1 && choice <= 4) {
            print_str("Enter first number: ");
            num1 = read_int();
            print_str("Enter second number: ");
            num2 = read_int();
            
            if (choice == 1) {
                result = add(num1, num2);
                print_str("Result: ");
                print_int(result);
                print_char('\n');
            } else if (choice == 2) {
                result = subtract(num1, num2);
                print_str("Result: ");
                print_int(result);
                print_char('\n');
            } else if (choice == 3) {
                result = multiply(num1, num2);
                print_str("Result: ");
                print_int(result);
                print_char('\n');
            } else if (choice == 4) {
                result = divide(num1, num2);
                if (num2 != 0) {
                    print_str("Result: ");
                    print_int(result);
                    print_char('\n');
                }
            }
        } else {
            print_str("Invalid choice!\n");
        }
    }
    
    exit(0);
}
\end{lstlisting}

% ============================================================================
% PART V: COMPILER IMPLEMENTATION
% ============================================================================
\part{Compiler Implementation}

% ============================================================================
% CHAPTER 28: COMPILER ARCHITECTURE
% ============================================================================
\chapter{Compiler Architecture}

\section{Overview}

The Algo compiler translates high-level Algo source code into A1M assembly language. It follows a traditional multi-pass compiler design with clear separation of concerns.

\subsection{Compilation Pipeline}

\begin{enumerate}
    \item \textbf{Lexical Analysis (Lexer):} Reads source code and produces stream of tokens
    \item \textbf{Syntax Analysis (Parser):} Builds Abstract Syntax Tree (AST) from tokens
    \item \textbf{Semantic Analysis:} Validates types, scopes, and semantic rules
    \item \textbf{Code Generation:} Converts AST to A1M assembly code
    \item \textbf{Assembly:} Assembles to binary executable (separate tool)
\end{enumerate}

\begin{figure}[H]
\centering
\begin{verbatim}
    hello.algo (Source Code)
           ↓
    ┌──────────────┐
    │    LEXER     │  Tokenization
    └──────────────┘
           ↓
      Token Stream
      [fun, identifier(main), lparen, rparen, ...]
           ↓
    ┌──────────────┐
    │    PARSER    │  Syntax Analysis
    └──────────────┘
           ↓
    Abstract Syntax Tree (AST)
      Program
        └── FunctionDecl (main)
              └── Block
                    └── CallExpr (print)
                          └── StringLiteral ("Hello")
           ↓
    ┌──────────────┐
    │   SEMANTIC   │  Type Checking
    └──────────────┘
           ↓
    Annotated AST (with type info)
           ↓
    ┌──────────────┐
    │   CODEGEN    │  Assembly Generation
    └──────────────┘
           ↓
    hello.asm (Assembly Code)
           ↓
    ┌──────────────┐
    │  ASSEMBLER   │  Binary Generation
    └──────────────┘
           ↓
    hello.bin (Executable)
\end{verbatim}
\caption{Complete Compilation Pipeline}
\end{figure}

\section{Compiler Implementation}

\subsection{Main Driver}

\begin{lstlisting}[style=c, caption={compiler.c - Main Compiler Driver}]
#include <stdio.h>
#include <stdlib.h>
#include <string.h>
#include "lexer.h"
#include "parser.h"
#include "semantic.h"
#include "codegen.h"

// Compiler options
typedef struct {
    char *input_file;
    char *output_file;
    int optimize;
    int verbose;
    int debug;
} CompilerOptions;

void print_usage(const char *program) {
    fprintf(stderr, "Usage: %s [options] input.algo\n", program);
    fprintf(stderr, "Options:\n");
    fprintf(stderr, "  -o <file>    Output file\n");
    fprintf(stderr, "  -O           Enable optimization\n");
    fprintf(stderr, "  -v           Verbose output\n");
    fprintf(stderr, "  -d           Debug mode\n");
}

int main(int argc, char **argv) {
    CompilerOptions opts = {0};
    
    // Parse command line arguments
    for (int i = 1; i < argc; i++) {
        if (strcmp(argv[i], "-o") == 0 && i + 1 < argc) {
            opts.output_file = argv[++i];
        } else if (strcmp(argv[i], "-O") == 0) {
            opts.optimize = 1;
        } else if (strcmp(argv[i], "-v") == 0) {
            opts.verbose = 1;
        } else if (strcmp(argv[i], "-d") == 0) {
            opts.debug = 1;
        } else if (argv[i][0] != '-') {
            opts.input_file = argv[i];
        } else {
            fprintf(stderr, "Unknown option: %s\n", argv[i]);
            print_usage(argv[0]);
            return 1;
        }
    }
    
    if (!opts.input_file) {
        fprintf(stderr, "Error: No input file specified\n");
        print_usage(argv[0]);
        return 1;
    }
    
    // Default output file
    if (!opts.output_file) {
        opts.output_file = "output.asm";
    }
    
    if (opts.verbose) {
        printf("Compiling: %s -> %s\n", 
               opts.input_file, opts.output_file);
    }
    
    // Phase 1: Lexical Analysis
    if (opts.verbose) printf("Phase 1: Lexical Analysis...\n");
    
    FILE *input = fopen(opts.input_file, "r");
    if (!input) {
        fprintf(stderr, "Error: Cannot open file '%s'\n", 
                opts.input_file);
        return 1;
    }
    
    Lexer *lexer = lexer_create(input);
    TokenList *tokens = lexer_tokenize(lexer);
    fclose(input);
    
    if (!tokens || lexer->error_count > 0) {
        fprintf(stderr, "Lexical analysis failed with %d errors\n", 
                lexer->error_count);
        return 1;
    }
    
    if (opts.debug) {
        printf("Tokens (%d):\n", tokens->count);
        for (int i = 0; i < tokens->count; i++) {
            token_print(&tokens->tokens[i]);
        }
    }
    
    // Phase 2: Syntax Analysis
    if (opts.verbose) printf("Phase 2: Syntax Analysis...\n");
    
    Parser *parser = parser_create(tokens);
    ASTNode *ast = parser_parse(parser);
    
    if (!ast || parser->error_count > 0) {
        fprintf(stderr, "Syntax analysis failed with %d errors\n", 
                parser->error_count);
        return 1;
    }
    
    if (opts.debug) {
        printf("AST:\n");
        ast_print(ast, 0);
    }
    
    // Phase 3: Semantic Analysis
    if (opts.verbose) printf("Phase 3: Semantic Analysis...\n");
    
    SemanticAnalyzer *analyzer = semantic_create();
    if (!semantic_analyze(analyzer, ast)) {
        fprintf(stderr, "Semantic analysis failed with %d errors\n", 
                analyzer->error_count);
        return 1;
    }
    
    if (opts.debug) {
        printf("Symbol Table:\n");
        symbol_table_print(analyzer->global_scope);
    }
    
    // Phase 4: Code Generation
    if (opts.verbose) printf("Phase 4: Code Generation...\n");
    
    FILE *output = fopen(opts.output_file, "w");
    if (!output) {
        fprintf(stderr, "Error: Cannot create output file '%s'\n", 
                opts.output_file);
        return 1;
    }
    
    CodeGenerator *codegen = codegen_create(output);
    codegen->optimize = opts.optimize;
    
    if (!codegen_generate(codegen, ast)) {
        fprintf(stderr, "Code generation failed\n");
        fclose(output);
        return 1;
    }
    
    fclose(output);
    
    // Success
    if (opts.verbose) {
        printf("Compilation successful!\n");
        printf("  Input:  %s\n", opts.input_file);
        printf("  Output: %s\n", opts.output_file);
        printf("  Tokens: %d\n", tokens->count);
        printf("  Lines:  %d\n", lexer->line);
    }
    
    // Cleanup
    lexer_free(lexer);
    token_list_free(tokens);
    parser_free(parser);
    semantic_free(analyzer);
    codegen_free(codegen);
    ast_free(ast);
    
    return 0;
}
\end{lstlisting}

\section{Data Structures}

\subsection{Token}

\begin{lstlisting}[style=c, caption={Token Structure}]
typedef enum {
    // Keywords
    TOKEN_FUN,      TOKEN_RETURN,   TOKEN_IF,
    TOKEN_ELSE,     TOKEN_WHILE,    TOKEN_FOR,
    TOKEN_INT,      TOKEN_BOOL,     TOKEN_CHAR,
    TOKEN_LONG,     TOKEN_TRUE,     TOKEN_FALSE,
    
    // Literals
    TOKEN_INTEGER_LITERAL,
    TOKEN_STRING_LITERAL,
    TOKEN_CHAR_LITERAL,
    
    // Identifiers
    TOKEN_IDENTIFIER,
    
    // Operators (Arithmetic)
    TOKEN_PLUS,     TOKEN_MINUS,    TOKEN_STAR,
    TOKEN_SLASH,    TOKEN_PERCENT,
    
    // Operators (Comparison)
    TOKEN_EQ,       TOKEN_NE,       TOKEN_LT,
    TOKEN_GT,       TOKEN_LE,       TOKEN_GE,
    
    // Operators (Logical)
    TOKEN_AND,      TOKEN_OR,       TOKEN_NOT,
    
    // Operators (Bitwise)
    TOKEN_BITAND,   TOKEN_BITOR,    TOKEN_BITXOR,
    TOKEN_BITNOT,   TOKEN_SHL,      TOKEN_SHR,
    
    // Assignment
    TOKEN_ASSIGN,
    
    // Delimiters
    TOKEN_LPAREN,   TOKEN_RPAREN,   TOKEN_LBRACE,
    TOKEN_RBRACE,   TOKEN_LBRACKET, TOKEN_RBRACKET,
    TOKEN_SEMICOLON,TOKEN_COMMA,    TOKEN_ARROW,
    
    // Special
    TOKEN_EOF,
    TOKEN_ERROR
} TokenType;

typedef struct {
    TokenType type;
    char *lexeme;        // Actual text
    int line;            // Line number
    int column;          // Column number
    
    // Value for literals
    union {
        int int_value;
        char *str_value;
        char char_value;
    } value;
} Token;

typedef struct {
    Token *tokens;
    int count;
    int capacity;
} TokenList;
    TOKEN_LBRACE,
    TOKEN_RBRACE,
    TOKEN_LBRACKET,
    TOKEN_RBRACKET,
    TOKEN_SEMICOLON,
    TOKEN_COMMA,
    TOKEN_ARROW,
    
    // Special
    TOKEN_EOF,
    TOKEN_ERROR
} TokenType;

typedef struct {
    TokenType type;
    char *value;
    int line;
    int column;
} Token;
\end{lstlisting}

\section{Lexer Implementation}

\begin{lstlisting}[style=c, caption={Lexer Core}]
typedef struct {
    char *source;
    size_t position;
    size_t length;
    int line;
    int column;
} Lexer;

void lexer_init(Lexer *lexer, const char *source) {
    lexer->source = strdup(source);
    lexer->position = 0;
    lexer->length = strlen(source);
    lexer->line = 1;
    lexer->column = 1;
}

char lexer_peek(Lexer *lexer) {
    if (lexer->position >= lexer->length) {
        return '\0';
    }
    return lexer->source[lexer->position];
}

char lexer_advance(Lexer *lexer) {
    char c = lexer_peek(lexer);
    lexer->position++;
    
    if (c == '\n') {
        lexer->line++;
        lexer->column = 1;
    } else {
        lexer->column++;
    }
    
    return c;
}

void lexer_skip_whitespace(Lexer *lexer) {
    while (isspace(lexer_peek(lexer))) {
        lexer_advance(lexer);
    }
}

void lexer_skip_comment(Lexer *lexer) {
    if (lexer_peek(lexer) == '/' && 
        lexer->source[lexer->position + 1] == '/') {
        // Single-line comment
        while (lexer_peek(lexer) != '\n' && 
               lexer_peek(lexer) != '\0') {
            lexer_advance(lexer);
        }
    } else if (lexer_peek(lexer) == '/' && 
               lexer->source[lexer->position + 1] == '*') {
        // Multi-line comment
        lexer_advance(lexer); // /
        lexer_advance(lexer); // *
        
        while (!(lexer_peek(lexer) == '*' && 
                lexer->source[lexer->position + 1] == '/')) {
            if (lexer_peek(lexer) == '\0') break;
            lexer_advance(lexer);
        }
        
        lexer_advance(lexer); // *
        lexer_advance(lexer); // /
    }
}

Token *lexer_next_token(Lexer *lexer) {
    lexer_skip_whitespace(lexer);
    lexer_skip_comment(lexer);
    lexer_skip_whitespace(lexer);
    
    char c = lexer_peek(lexer);
    
    if (c == '\0') {
        return create_token(TOKEN_EOF, "", 
                          lexer->line, lexer->column);
    }
    
    // Identifiers and keywords
    if (isalpha(c) || c == '_') {
        return lexer_identifier(lexer);
    }
    
    // Numbers
    if (isdigit(c)) {
        return lexer_number(lexer);
    }
    
    // Strings
    if (c == '"') {
        return lexer_string(lexer);
    }
    
    // Operators and delimiters
    return lexer_operator(lexer);
}
\end{lstlisting}

% ============================================================================
% CHAPTER 26: SYNTAX ANALYSIS
% ============================================================================
\chapter{Syntax Analysis}

\section{Abstract Syntax Tree}

\begin{lstlisting}[style=c, caption={AST Node Types}]
typedef enum {
    AST_PROGRAM,
    AST_FUNCTION,
    AST_PARAMETER,
    AST_BLOCK,
    AST_RETURN,
    AST_IF,
    AST_WHILE,
    AST_FOR,
    AST_VAR_DECL,
    AST_ASSIGN,
    AST_BINARY_OP,
    AST_UNARY_OP,
    AST_CALL,
    AST_ARRAY_ACCESS,
    AST_IDENTIFIER,
    AST_LITERAL
} ASTNodeType;

typedef struct ASTNode {
    ASTNodeType type;
    DataType data_type;
    
    union {
        // Function
        struct {
            char *name;
            struct ASTNode **parameters;
            int param_count;
            struct ASTNode *body;
            DataType return_type;
        } function;
        
        // Binary operation
        struct {
            char *op;
            struct ASTNode *left;
            struct ASTNode *right;
        } binary_op;
        
        // Variable declaration
        struct {
            char *name;
            DataType type;
            struct ASTNode *initializer;
        } var_decl;
        
        // Literal
        struct {
            union {
                int64_t int_value;
                char *str_value;
                char char_value;
                bool bool_value;
            };
        } literal;
        
        // ... other node types
    };
    
    struct ASTNode **children;
    int child_count;
} ASTNode;
\end{lstlisting}

\section{Parser Implementation}

\begin{lstlisting}[style=c, caption={Recursive Descent Parser}]
typedef struct {
    Token **tokens;
    int position;
    int count;
} Parser;

void parser_init(Parser *p, Token **tokens, int count) {
    p->tokens = tokens;
    p->position = 0;
    p->count = count;
}

Token *parser_peek(Parser *p) {
    if (p->position >= p->count) return NULL;
    return p->tokens[p->position];
}

Token *parser_advance(Parser *p) {
    Token *token = parser_peek(p);
    if (token) p->position++;
    return token;
}

bool parser_expect(Parser *p, TokenType type) {
    Token *token = parser_peek(p);
    if (!token || token->type != type) {
        fprintf(stderr, "Parse error at line %d: expected %d\n",
                token ? token->line : 0, type);
        return false;
    }
    parser_advance(p);
    return true;
}

// Parse program: function*
ASTNode *parse_program(Parser *p) {
    ASTNode *program = create_node(AST_PROGRAM);
    
    while (parser_peek(p)->type != TOKEN_EOF) {
        ASTNode *func = parse_function(p);
        if (!func) return NULL;
        add_child(program, func);
    }
    
    return program;
}

// Parse function: 'fun' ID '(' params ')' '->' type block
ASTNode *parse_function(Parser *p) {
    if (!parser_expect(p, TOKEN_FUN)) return NULL;
    
    Token *name = parser_peek(p);
    if (!parser_expect(p, TOKEN_IDENTIFIER)) return NULL;
    
    if (!parser_expect(p, TOKEN_LPAREN)) return NULL;
    
    // Parse parameters
    ASTNode **params = NULL;
    int param_count = 0;
    
    while (parser_peek(p)->type != TOKEN_RPAREN) {
        ASTNode *param = parse_parameter(p);
        if (!param) return NULL;
        
        params = realloc(params, 
                        sizeof(ASTNode*) * (param_count + 1));
        params[param_count++] = param;
        
        if (parser_peek(p)->type == TOKEN_COMMA) {
            parser_advance(p);
        }
    }
    
    if (!parser_expect(p, TOKEN_RPAREN)) return NULL;
    
    // Parse return type
    DataType return_type = TYPE_VOID;
    if (parser_peek(p)->type == TOKEN_ARROW) {
        parser_advance(p);
        return_type = parse_type(p);
    }
    
    // Parse body
    ASTNode *body = parse_block(p);
    if (!body) return NULL;
    
    // Create function node
    ASTNode *func = create_function_node(name->value, 
                                         params, param_count,
                                         body, return_type);
    return func;
}

// Parse expression with precedence climbing
ASTNode *parse_expression(Parser *p, int min_prec) {
    ASTNode *left = parse_primary(p);
    if (!left) return NULL;
    
    while (true) {
        Token *op = parser_peek(p);
        if (!is_binary_op(op->type)) break;
        
        int prec = get_precedence(op->type);
        if (prec < min_prec) break;
        
        parser_advance(p);
        
        ASTNode *right = parse_expression(p, prec + 1);
        if (!right) return NULL;
        
        left = create_binary_node(op->value, left, right);
    }
    
    return left;
}
\end{lstlisting}

% ============================================================================
% CHAPTER 27: SEMANTIC ANALYSIS
% ============================================================================
\chapter{Semantic Analysis}

\section{Symbol Table}

\begin{lstlisting}[style=c, caption={Symbol Table Implementation}]
typedef struct Symbol {
    char *name;
    DataType type;
    int offset;          // Stack offset for locals
    bool is_global;
    bool is_parameter;
    struct Symbol *next;
} Symbol;

typedef struct SymbolTable {
    Symbol *symbols;
    struct SymbolTable *parent;
    int stack_offset;
} SymbolTable;

SymbolTable *symbol_table_create(SymbolTable *parent) {
    SymbolTable *table = malloc(sizeof(SymbolTable));
    table->symbols = NULL;
    table->parent = parent;
    table->stack_offset = parent ? parent->stack_offset : 0;
    return table;
}

Symbol *symbol_table_lookup(SymbolTable *table, 
                            const char *name) {
    // Search current scope
    Symbol *sym = table->symbols;
    while (sym) {
        if (strcmp(sym->name, name) == 0) {
            return sym;
        }
        sym = sym->next;
    }
    
    // Search parent scopes
    if (table->parent) {
        return symbol_table_lookup(table->parent, name);
    }
    
    return NULL;
}

bool symbol_table_insert(SymbolTable *table, 
                        const char *name, DataType type) {
    // Check if already exists in current scope
    Symbol *existing = table->symbols;
    while (existing) {
        if (strcmp(existing->name, name) == 0) {
            return false; // Duplicate
        }
        existing = existing->next;
    }
    
    // Create new symbol
    Symbol *sym = malloc(sizeof(Symbol));
    sym->name = strdup(name);
    sym->type = type;
    sym->offset = table->stack_offset;
    sym->is_global = (table->parent == NULL);
    sym->next = table->symbols;
    
    table->symbols = sym;
    table->stack_offset += get_type_size(type);
    
    return true;
}
\end{lstlisting}

\section{Type Checking}

\begin{lstlisting}[style=c, caption={Type Checker}]
bool semantic_analyze_node(ASTNode *node, 
                          SymbolTable *table) {
    switch (node->type) {
        case AST_FUNCTION:
            return analyze_function(node, table);
            
        case AST_VAR_DECL:
            if (!symbol_table_insert(table, 
                                    node->var_decl.name,
                                    node->var_decl.type)) {
                fprintf(stderr, "Duplicate variable: %s\n",
                       node->var_decl.name);
                return false;
            }
            
            if (node->var_decl.initializer) {
                if (!semantic_analyze_node(
                        node->var_decl.initializer, table)) {
                    return false;
                }
                
                // Type compatibility check
                if (node->var_decl.type != 
                    node->var_decl.initializer->data_type) {
                    fprintf(stderr, "Type mismatch\n");
                    return false;
                }
            }
            return true;
            
        case AST_BINARY_OP:
            if (!semantic_analyze_node(node->binary_op.left, 
                                      table)) {
                return false;
            }
            if (!semantic_analyze_node(node->binary_op.right, 
                                      table)) {
                return false;
            }
            
            // Determine result type
            DataType left_type = node->binary_op.left->data_type;
            DataType right_type = node->binary_op.right->data_type;
            
            if (left_type != right_type) {
                fprintf(stderr, "Type mismatch in binary op\n");
                return false;
            }
            
            node->data_type = left_type;
            return true;
            
        case AST_IDENTIFIER:
            {
                Symbol *sym = symbol_table_lookup(table, 
                                                 node->literal.str_value);
                if (!sym) {
                    fprintf(stderr, "Undefined variable: %s\n",
                           node->literal.str_value);
                    return false;
                }
                node->data_type = sym->type;
            }
            return true;
            
        default:
            // Analyze children
            for (int i = 0; i < node->child_count; i++) {
                if (!semantic_analyze_node(node->children[i], 
                                          table)) {
                    return false;
                }
            }
            return true;
    }
}
\end{lstlisting}

% ============================================================================
% CHAPTER 28: CODE GENERATION
% ============================================================================
\chapter{Code Generation}

\section{Code Generator Architecture}

\begin{lstlisting}[style=c, caption={Code Generator}]
typedef struct {
    FILE *output;
    int label_counter;
    int temp_counter;
    SymbolTable *current_scope;
} CodeGen;

void codegen_init(CodeGen *cg, FILE *output) {
    cg->output = output;
    cg->label_counter = 0;
    cg->temp_counter = 10; // Start from R10
}

void emit(CodeGen *cg, const char *fmt, ...) {
    va_list args;
    va_start(args, fmt);
    vfprintf(cg->output, fmt, args);
    va_end(args);
    fprintf(cg->output, "\n");
}

int new_label(CodeGen *cg) {
    return cg->label_counter++;
}

int alloc_temp_reg(CodeGen *cg) {
    return cg->temp_counter++;
}

void generate_function(CodeGen *cg, ASTNode *func) {
    emit(cg, "\n; Function: %s", func->function.name);
    emit(cg, "%s:", func->function.name);
    
    // Function prologue
    emit(cg, "    PUSH R31");  // Save return address
    
    // Allocate stack space for locals
    int local_size = calculate_local_size(func);
    if (local_size > 0) {
        emit(cg, "    SUBI R3, R3, %d", local_size);
    }
    
    // Generate function body
    generate_statement(cg, func->function.body);
    
    // Function epilogue
    if (local_size > 0) {
        emit(cg, "    ADDI R3, R3, %d", local_size);
    }
    emit(cg, "    POP R31");
    emit(cg, "    RET");
}

void generate_expression(CodeGen *cg, ASTNode *expr, 
                        int dest_reg) {
    switch (expr->type) {
        case AST_LITERAL:
            emit(cg, "    MOVI R%d, %lld", dest_reg,
                 expr->literal.int_value);
            break;
            
        case AST_IDENTIFIER:
            {
                Symbol *sym = symbol_table_lookup(
                    cg->current_scope, 
                    expr->literal.str_value);
                emit(cg, "    LOAD R%d, %d(R3)", 
                     dest_reg, sym->offset);
            }
            break;
            
        case AST_BINARY_OP:
            {
                int left_reg = alloc_temp_reg(cg);
                int right_reg = alloc_temp_reg(cg);
                
                generate_expression(cg, expr->binary_op.left, 
                                   left_reg);
                generate_expression(cg, expr->binary_op.right, 
                                   right_reg);
                
                char *op = expr->binary_op.op;
                if (strcmp(op, "+") == 0) {
                    emit(cg, "    ADD R%d, R%d, R%d", 
                         dest_reg, left_reg, right_reg);
                } else if (strcmp(op, "-") == 0) {
                    emit(cg, "    SUB R%d, R%d, R%d", 
                         dest_reg, left_reg, right_reg);
                } else if (strcmp(op, "*") == 0) {
                    emit(cg, "    MUL R%d, R%d, R%d", 
                         dest_reg, left_reg, right_reg);
                } else if (strcmp(op, "/") == 0) {
                    emit(cg, "    DIV R%d, R%d, R%d", 
                         dest_reg, left_reg, right_reg);
                }
                
                cg->temp_counter -= 2;
            }
            break;
            
        case AST_CALL:
            generate_call(cg, expr, dest_reg);
            break;
    }
}

void generate_if(CodeGen *cg, ASTNode *if_node) {
    int else_label = new_label(cg);
    int end_label = new_label(cg);
    
    // Evaluate condition
    int cond_reg = alloc_temp_reg(cg);
    generate_expression(cg, if_node->children[0], cond_reg);
    
    // Branch to else if condition is false
    emit(cg, "    CMPI R%d, 0", cond_reg);
    emit(cg, "    JZ L%d", else_label);
    
    // Then block
    generate_statement(cg, if_node->children[1]);
    emit(cg, "    JMP L%d", end_label);
    
    // Else block
    emit(cg, "L%d:", else_label);
    if (if_node->child_count > 2) {
        generate_statement(cg, if_node->children[2]);
    }
    
    emit(cg, "L%d:", end_label);
    cg->temp_counter--;
}
\end{lstlisting}

% ============================================================================
% PART VI: COMPILATION PROCESS
% ============================================================================
\part{Compilation Process}

% ============================================================================
% CHAPTER 29: BUILD SYSTEM
% ============================================================================
\chapter{Build System and Toolchain}

\section{Complete Build Pipeline}

The A1M toolchain consists of three main tools that work together:

\begin{enumerate}
    \item \textbf{Compiler (algo):} Algo source → Assembly
    \item \textbf{Assembler (a1masm):} Assembly → Binary
    \item \textbf{Emulator (a1memu):} Binary execution
\end{enumerate}

\subsection{Build Script}

\begin{lstlisting}[style=c, caption={build.sh - Complete Build Script}]
#!/bin/bash

set -e  # Exit on error

ALGO_FILE=$1
BASE_NAME=$(basename "$ALGO_FILE" .algo)

echo "=== A1M Build System ==="
echo "Source: $ALGO_FILE"

# Step 1: Compile Algo to Assembly
echo "[1/3] Compiling..."
algo "$ALGO_FILE" "${BASE_NAME}.asm"

# Step 2: Assemble to Binary
echo "[2/3] Assembling..."
a1masm "${BASE_NAME}.asm" "${BASE_NAME}.bin"

# Step 3: Run
echo "[3/3] Running..."
a1memu "${BASE_NAME}.bin"

echo "=== Build Complete ==="
\end{lstlisting}

\section{Makefile}

\begin{lstlisting}[caption={Makefile for A1M Projects}]
CC = gcc
CFLAGS = -Wall -Wextra -O2

COMPILER_SRC = compiler/src
ASSEMBLER_SRC = assembler/src
EMULATOR_SRC = emulator/src

.PHONY: all clean install

all: algo a1masm a1memu

# Compiler
algo: $(COMPILER_SRC)/*.c
	$(CC) $(CFLAGS) -o algo $(COMPILER_SRC)/*.c

# Assembler
a1masm: $(ASSEMBLER_SRC)/*.c
	$(CC) $(CFLAGS) -o a1masm $(ASSEMBLER_SRC)/*.c

# Emulator
a1memu: $(EMULATOR_SRC)/*.c
	$(CC) $(CFLAGS) -o a1memu $(EMULATOR_SRC)/*.c

install:
	install -m 755 algo /usr/local/bin/
	install -m 755 a1masm /usr/local/bin/
	install -m 755 a1memu /usr/local/bin/

clean:
	rm -f algo a1masm a1memu
	rm -f *.asm *.bin
\end{lstlisting}

% ============================================================================
% CHAPTER 30: OPTIMIZATION TECHNIQUES
% ============================================================================
\chapter{Optimization Techniques}

\section{Constant Folding}

Evaluate constant expressions at compile time:

\begin{lstlisting}[style=algo]
int x = 3 + 5 * 2;  // Optimized to: int x = 13;
\end{lstlisting}

\section{Dead Code Elimination}

Remove unreachable code:

\begin{lstlisting}[style=algo]
if (false) {
    // This code is eliminated
    print("Never executed");
}
\end{lstlisting}

\section{Register Allocation}

Efficient use of the 32 available registers minimizes memory access.

% ============================================================================
% PART VII: PRACTICAL EXAMPLES
% ============================================================================
\part{Practical Examples and Applications}

% ============================================================================
% CHAPTER 31: EXAMPLE PROGRAMS
% ============================================================================
\chapter{Complete Example Programs}

\section{Fibonacci Calculator}

\begin{lstlisting}[style=algo, caption={Fibonacci Program}]
fun fibonacci(int n) -> int {
    if (n <= 1) {
        return n;
    }
    return fibonacci(n - 1) + fibonacci(n - 2);
}

fun main() {
    print("Fibonacci sequence:\n");
    for (int i = 0; i < 15; i = i + 1) {
        print_int(fibonacci(i));
        print(" ");
    }
    print("\n");
}
\end{lstlisting}

\section{Sorting Algorithms}

\begin{lstlisting}[style=algo, caption={Quick Sort}]
fun partition(int arr[], int low, int high) -> int {
    int pivot = arr[high];
    int i = low - 1;
    
    for (int j = low; j < high; j = j + 1) {
        if (arr[j] < pivot) {
            i = i + 1;
            int temp = arr[i];
            arr[i] = arr[j];
            arr[j] = temp;
        }
    }
    
    int temp = arr[i + 1];
    arr[i + 1] = arr[high];
    arr[high] = temp;
    
    return i + 1;
}

fun quicksort(int arr[], int low, int high) {
    if (low < high) {
        int pi = partition(arr, low, high);
        quicksort(arr, low, pi - 1);
        quicksort(arr, pi + 1, high);
    }
}
\end{lstlisting}

% ============================================================================
% CHAPTER 32: DEBUGGING TECHNIQUES
% ============================================================================
\chapter{Debugging and Testing}

\section{Using the Emulator Debugger}

\begin{verbatim}
$ a1memu -d program.bin

=== A1M Emulator (Debug Mode) ===
PC: 0x00000000  Instruction: MOVI R10, 42

Commands:
  s - Step one instruction
  c - Continue execution
  r - Show registers
  m <addr> - Show memory
  q - Quit
\end{verbatim}

\section{Common Debugging Techniques}

\begin{itemize}
    \item \textbf{Print debugging:} Use print statements
    \item \textbf{Register inspection:} Check register values
    \item \textbf{Memory dumps:} Examine memory contents
    \item \textbf{Single-stepping:} Execute one instruction at a time
\end{itemize}

% ============================================================================
% CHAPTER 33: PERFORMANCE ANALYSIS
% ============================================================================
\chapter{Performance Analysis}

\section{Benchmark Results}

\begin{table}[H]
\centering
\caption{A1M Performance Benchmarks}
\begin{tabular}{|l|r|r|}
\hline
\textbf{Benchmark} & \textbf{Cycles} & \textbf{Instructions} \\
\hline
Fibonacci(20) & 45,123 & 38,941 \\
Bubble Sort (100) & 125,678 & 98,234 \\
Quick Sort (100) & 23,456 & 19,876 \\
Matrix Multiply (10×10) & 67,890 & 54,321 \\
\hline
\end{tabular}
\end{table}

% ============================================================================
% PART VIII: TESTING AND FUTURE WORK
% ============================================================================
\part{Testing and Future Enhancements}

% ============================================================================
% CHAPTER 34: TEST SUITE
% ============================================================================
\chapter{Comprehensive Test Suite}

\section{Unit Tests}

Testing individual compiler components:

\begin{lstlisting}[style=c, caption={Lexer Tests}]
void test_lexer_integer() {
    Lexer lexer;
    lexer_init(&lexer, "42");
    
    Token *token = lexer_next_token(&lexer);
    assert(token->type == TOKEN_INTEGER);
    assert(strcmp(token->value, "42") == 0);
}

void test_lexer_identifier() {
    Lexer lexer;
    lexer_init(&lexer, "variable");
    
    Token *token = lexer_next_token(&lexer);
    assert(token->type == TOKEN_IDENTIFIER);
}
\end{lstlisting}

\section{Integration Tests}

End-to-end testing of complete programs:

\begin{lstlisting}[style=algo, caption={Test Program}]
// test_arithmetic.algo
fun main() {
    int a = 10;
    int b = 20;
    int c = a + b;
    
    if (c == 30) {
        print("PASS\n");
    } else {
        print("FAIL\n");
    }
}
\end{lstlisting}

% ============================================================================
% CHAPTER 35: FUTURE ENHANCEMENTS
% ============================================================================
\chapter{Future Enhancements}

\section{Planned Features}

\subsection{Language Extensions}

\begin{itemize}
    \item \textbf{Pointers:} Direct memory manipulation
    \item \textbf{Structures:} User-defined types
    \item \textbf{Floating-point:} FP arithmetic support
    \item \textbf{Strings:} First-class string type
\end{itemize}

\subsection{Hardware Improvements}

\begin{itemize}
    \item \textbf{Cache:} L1/L2 cache hierarchy
    \item \textbf{Branch Prediction:} Dynamic prediction
    \item \textbf{Out-of-Order Execution:} Superscalar pipeline
    \item \textbf{SIMD:} Vector operations
\end{itemize}

\subsection{Compiler Optimizations}

\begin{itemize}
    \item \textbf{Loop Unrolling:} Reduce loop overhead
    \item \textbf{Inlining:} Eliminate function call overhead
    \item \textbf{Common Subexpression Elimination}
    \item \textbf{Register Allocation:} Graph coloring algorithm
\end{itemize}

\section{Conclusion}

The A1M architecture demonstrates a complete computer system from hardware to high-level programming language. This thesis has presented:

\begin{enumerate}
    \item A clean 64-bit RISC architecture with 32 registers
    \item Complete Verilog HDL implementation
    \item Two-pass assembler for machine code generation
    \item Cycle-accurate emulator for testing
    \item High-level Algo programming language
    \item Full compiler with lexer, parser, and code generator
    \item VS Code integration for modern development
\end{enumerate}

The A1M system serves both as a functional computing platform and an educational tool for understanding computer architecture and compiler design principles.

% ============================================================================
% APPENDICES
% ============================================================================
\appendix

% ============================================================================
% APPENDIX A: INSTRUCTION REFERENCE
% ============================================================================
\chapter{Complete Instruction Reference}

\section{Instruction Summary}

\begin{table}[H]
\centering
\caption{Complete A1M Instruction Set}
\small
\begin{tabularx}{\textwidth}{|l|l|l|X|}
\hline
\textbf{Opcode} & \textbf{Mnemonic} & \textbf{Format} & \textbf{Description} \\
\hline
0x00 & NOP & - & No operation \\
0x01 & HALT & - & Halt execution \\
0x02 & ADD & R & Add registers \\
0x03 & ADDI & I & Add immediate \\
0x04 & SUB & R & Subtract registers \\
0x05 & SUBI & I & Subtract immediate \\
0x06 & MUL & R & Multiply registers \\
0x07 & MULI & I & Multiply immediate \\
0x08 & DIV & R & Divide registers \\
0x09 & DIVI & I & Divide immediate \\
0x0A & MOD & R & Modulo registers \\
0x0B & MODI & I & Modulo immediate \\
0x0C & AND & R & Bitwise AND \\
0x0D & ANDI & I & AND immediate \\
0x0E & OR & R & Bitwise OR \\
0x0F & ORI & I & OR immediate \\
0x10 & XOR & R & Bitwise XOR \\
0x11 & XORI & I & XOR immediate \\
0x12 & SHL & R & Shift left \\
0x13 & SHLI & I & Shift left immediate \\
0x14 & SHR & R & Shift right logical \\
0x15 & SHRI & I & Shift right immediate \\
0x16 & SAR & R & Shift right arithmetic \\
0x17 & SARI & I & SAR immediate \\
0x20 & LOAD & M & Load 64-bit word \\
0x21 & STORE & M & Store 64-bit word \\
0x22 & LOADB & M & Load byte \\
0x23 & STOREB & M & Store byte \\
0x30 & JMP & J & Unconditional jump \\
0x31 & JZ & J & Jump if zero \\
0x32 & JNZ & J & Jump if not zero \\
0x33 & JGT & J & Jump if greater \\
0x34 & JLT & J & Jump if less \\
0x35 & JGE & J & Jump if greater/equal \\
0x36 & JLE & J & Jump if less/equal \\
0x39 & CALL & J & Function call \\
0x3A & RET & - & Return from function \\
0x40 & CMP & R & Compare registers \\
0x41 & CMPI & I & Compare immediate \\
0x42 & PUSH & R & Push to stack \\
0x43 & POP & R & Pop from stack \\
0x50 & SYSCALL & I & System call \\
\hline
\end{tabularx}
\end{table}

% ============================================================================
% APPENDIX B: REGISTER CONVENTIONS
% ============================================================================
\chapter{Register Usage Conventions}

\begin{table}[H]
\centering
\caption{A1M Calling Convention}
\begin{tabularx}{\textwidth}{|l|l|l|X|}
\hline
\textbf{Register} & \textbf{Name} & \textbf{Preserved} & \textbf{Purpose} \\
\hline
R0 & zero & N/A & Always zero \\
R1 & ret/arg1 & No & Return value / Syscall arg 1 \\
R2 & arg2 & No & Syscall argument 2 \\
R3 & sp & Yes & Stack pointer \\
R4-R11 & a0-a7 & No & Function arguments \\
R12-R30 & t0-t18 & No & Temporaries \\
R31 & ra & No & Return address \\
\hline
\end{tabularx}
\end{table}

% ============================================================================
% APPENDIX C: SYSTEM CALL REFERENCE
% ============================================================================
\chapter{System Call Reference}

\begin{table}[H]
\centering
\caption{System Call Numbers and Arguments}
\small
\begin{tabularx}{\textwidth}{|l|l|X|}
\hline
\textbf{Number} & \textbf{Name} & \textbf{Arguments (R2, R3, R4)} \\
\hline
0 & exit & exit\_code \\
1 & write & fd, buffer, length \\
2 & read & fd, buffer, length \\
3 & open & path, flags \\
4 & close & fd \\
5 & malloc & size → returns address in R1 \\
6 & free & address \\
10 & print\_int & value \\
11 & print\_str & string\_address \\
12 & read\_int & → returns value in R1 \\
13 & read\_str & buffer, max\_length \\
\hline
\end{tabularx}
\end{table}

% ============================================================================
% APPENDIX D: ALGO LANGUAGE GRAMMAR
% ============================================================================
\chapter{Algo Language Grammar}

\section{EBNF Grammar}

\begin{verbatim}
program     ::= function*

function    ::= 'fun' IDENTIFIER '(' parameters ')' 
                ('->' type)? block

parameters  ::= parameter (',' parameter)*
parameter   ::= type IDENTIFIER ('[' ']')?

block       ::= '{' statement* '}'

statement   ::= var_decl
              | assign
              | if_stmt
              | while_stmt
              | for_stmt
              | return_stmt
              | expr_stmt

var_decl    ::= type IDENTIFIER ('[' INTEGER ']')? 
                ('=' expression)? ';'

assign      ::= IDENTIFIER ('['expression']')? '=' expression ';'

if_stmt     ::= 'if' '(' expression ')' block ('else' block)?

while_stmt  ::= 'while' '(' expression ')' block

for_stmt    ::= 'for' '(' var_decl expression ';' 
                assign ')' block

return_stmt ::= 'return' expression? ';'

expression  ::= logic_or

logic_or    ::= logic_and ('||' logic_and)*
logic_and   ::= equality ('&&' equality)*
equality    ::= comparison (('=='|'!=') comparison)*
comparison  ::= term (('<'|'>'|'<='|'>=') term)*
term        ::= factor (('+'|'-') factor)*
factor      ::= unary (('*'|'/'|'%') unary)*
unary       ::= ('!'|'-') unary | primary

primary     ::= INTEGER | STRING | CHAR | 'true' | 'false'
              | IDENTIFIER
              | IDENTIFIER '(' arguments ')'
              | IDENTIFIER '[' expression ']'
              | '(' expression ')'

type        ::= 'int' | 'long' | 'char' | 'bool'
\end{verbatim}

% ============================================================================
% APPENDICES
% ============================================================================
\appendix

% ============================================================================
% APPENDIX A: COMPLETE INSTRUCTION REFERENCE
% ============================================================================
\chapter{Complete Instruction Reference}

\section{Instruction Set Summary}

This appendix provides a complete reference for all A1M instructions, including opcodes, formats, operands, and detailed descriptions.

\subsection{Instruction Format Legend}

\begin{itemize}
    \item \textbf{RD} - Destination register (5 bits, R0-R31)
    \item \textbf{RS1} - Source register 1 (5 bits, R0-R31)
    \item \textbf{RS2} - Source register 2 (5 bits, R0-R31)
    \item \textbf{IMM} - Immediate value (14 bits, signed)
    \item \textbf{ADDR} - Jump/branch address (24 bits)
    \item \textbf{OFFSET} - Memory offset (14 bits, signed)
\end{itemize}

\section{Complete Instruction List}

\begin{table}[H]
\centering
\caption{A1M Complete Instruction Set}
\label{tab:complete-instructions}
\tiny
\begin{tabularx}{\textwidth}{|l|l|l|X|l|}
\hline
\textbf{Opcode} & \textbf{Mnemonic} & \textbf{Format} & \textbf{Operation} & \textbf{Flags} \\
\hline
\multicolumn{5}{|c|}{\textbf{Control Instructions}} \\
\hline
0x00 & NOP & - & No operation & - \\
0x01 & HALT & - & Stop execution & - \\
\hline
\multicolumn{5}{|c|}{\textbf{Arithmetic Instructions}} \\
\hline
0x02 & ADD & R & RD = RS1 + RS2 & ZF, NF \\
0x03 & ADDI & I & RD = RS1 + IMM & ZF, NF \\
0x04 & SUB & R & RD = RS1 - RS2 & ZF, NF \\
0x05 & SUBI & I & RD = RS1 - IMM & ZF, NF \\
0x06 & MUL & R & RD = RS1 $\times$ RS2 & ZF, NF \\
0x07 & MULI & I & RD = RS1 $\times$ IMM & ZF, NF \\
0x08 & DIV & R & RD = RS1 / RS2 & ZF, NF \\
0x09 & DIVI & I & RD = RS1 / IMM & ZF, NF \\
0x0A & MOD & R & RD = RS1 \% RS2 & ZF, NF \\
0x0B & MODI & I & RD = RS1 \% IMM & ZF, NF \\
\hline
\multicolumn{5}{|c|}{\textbf{Logical Instructions}} \\
\hline
0x0C & AND & R & RD = RS1 \& RS2 & ZF, NF \\
0x0D & ANDI & I & RD = RS1 \& IMM & ZF, NF \\
0x0E & OR & R & RD = RS1 | RS2 & ZF, NF \\
0x0F & ORI & I & RD = RS1 | IMM & ZF, NF \\
0x10 & XOR & R & RD = RS1 $\oplus$ RS2 & ZF, NF \\
0x11 & XORI & I & RD = RS1 $\oplus$ IMM & ZF, NF \\
0x12 & SHL & R & RD = RS1 $<<$ RS2 & ZF, NF \\
0x13 & SHLI & I & RD = RS1 $<<$ IMM & ZF, NF \\
0x14 & SHR & R & RD = RS1 $>>$ RS2 (logical) & ZF, NF \\
0x15 & SHRI & I & RD = RS1 $>>$ IMM (logical) & ZF, NF \\
0x16 & SAR & R & RD = RS1 $>>$ RS2 (arithmetic) & ZF, NF \\
0x17 & SARI & I & RD = RS1 $>>$ IMM (arithmetic) & ZF, NF \\
0x18 & NOT & R & RD = $\sim$RS1 & ZF, NF \\
\hline
\multicolumn{5}{|c|}{\textbf{Memory Instructions}} \\
\hline
0x20 & LOAD & M & RD = Mem[RS1 + OFFSET] & - \\
0x21 & STORE & M & Mem[RS1 + OFFSET] = RD & - \\
0x22 & LOADB & M & RD = (byte)Mem[RS1 + OFFSET] & - \\
0x23 & STOREB & M & Mem[RS1 + OFFSET] = (byte)RD & - \\
\hline
\multicolumn{5}{|c|}{\textbf{Branch and Jump Instructions}} \\
\hline
0x30 & JMP & J & PC = ADDR & - \\
0x31 & JZ & J & if (ZF) PC = ADDR & - \\
0x32 & JNZ & J & if (!ZF) PC = ADDR & - \\
0x33 & JGT & J & if (!ZF \&\& !NF) PC = ADDR & - \\
0x34 & JLT & J & if (NF) PC = ADDR & - \\
0x35 & JGE & J & if (!NF) PC = ADDR & - \\
0x36 & JLE & J & if (ZF || NF) PC = ADDR & - \\
0x37 & JEQ & J & if (ZF) PC = ADDR & - \\
0x38 & JNE & J & if (!ZF) PC = ADDR & - \\
0x39 & CALL & J & R31 = PC + 4; PC = ADDR & - \\
0x3A & RET & - & PC = R31 & - \\
\hline
\multicolumn{5}{|c|}{\textbf{Comparison Instructions}} \\
\hline
0x40 & CMP & R & Set flags: RS1 - RS2 & ZF, NF, CF, OF \\
0x41 & CMPI & I & Set flags: RS1 - IMM & ZF, NF, CF, OF \\
\hline
\multicolumn{5}{|c|}{\textbf{Stack Instructions}} \\
\hline
0x42 & PUSH & R & SP -= 8; Mem[SP] = RS1 & - \\
0x43 & POP & R & RD = Mem[SP]; SP += 8 & - \\
\hline
\multicolumn{5}{|c|}{\textbf{System Instructions}} \\
\hline
0x50 & SYSCALL & I & System call (IMM = number) & - \\
\hline
\end{tabularx}
\end{table}

\section{Instruction Encoding Details}

\subsection{R-Type Format}

\begin{verbatim}
 31    24 23  19 18  14 13   9 8       0
┌────────┬──────┬──────┬──────┬─────────┐
│ Opcode │  RD  │ RS1  │ RS2  │ Unused  │
└────────┴──────┴──────┴──────┴─────────┘
  8 bits  5 bits 5 bits 5 bits  9 bits
\end{verbatim}

\textbf{Example:} \code{ADD R5, R10, R15}

\begin{itemize}
    \item Opcode: 0x02 (ADD)
    \item RD: 5 (00101)
    \item RS1: 10 (01010)
    \item RS2: 15 (01111)
    \item Encoding: \code{0x02 2A 78 00}
\end{itemize}

\subsection{I-Type Format}

\begin{verbatim}
 31    24 23  19 18  14 13           0
┌────────┬──────┬──────┬──────────────┐
│ Opcode │  RD  │ RS1  │  Immediate   │
└────────┴──────┴──────┴──────────────┘
  8 bits  5 bits 5 bits   14 bits
\end{verbatim}

\textbf{Example:} \code{ADDI R3, R4, 100}

\begin{itemize}
    \item Opcode: 0x03 (ADDI)
    \item RD: 3 (00011)
    \item RS1: 4 (00100)
    \item Immediate: 100 (0000000001100100)
    \item Encoding: \code{0x03 1A 06 64}
\end{itemize}

% ============================================================================
% APPENDIX B: REGISTER CONVENTIONS
% ============================================================================
\chapter{Register Usage Conventions}

\section{Register Assignment}

\begin{table}[H]
\centering
\caption{A1M Register Assignment and Usage}
\label{tab:register-assignment}
\begin{tabularx}{\textwidth}{|l|l|l|l|X|}
\hline
\textbf{Reg} & \textbf{ABI} & \textbf{Saved} & \textbf{Usage} & \textbf{Notes} \\
\hline
R0 & zero & N/A & Constant zero & Hardwired to 0, writes ignored \\
R1 & ret & No & Return value & Also syscall number \\
R2 & arg2 & No & Syscall arg 2 & Second syscall argument \\
R3 & sp & Yes & Stack pointer & Must be preserved by callee \\
R4 & a0 & No & Argument 0 & First function argument \\
R5 & a1 & No & Argument 1 & Second function argument \\
R6 & a2 & No & Argument 2 & Third function argument \\
R7 & a3 & No & Argument 3 & Fourth function argument \\
R8 & a4 & No & Argument 4 & Fifth function argument \\
R9 & a5 & No & Argument 5 & Sixth function argument \\
R10 & a6 & No & Argument 6 & Seventh function argument \\
R11 & a7 & No & Argument 7 & Eighth function argument \\
R12 & t0 & No & Temporary 0 & Caller-saved temporary \\
R13 & t1 & No & Temporary 1 & Caller-saved temporary \\
R14 & t2 & No & Temporary 2 & Caller-saved temporary \\
R15 & t3 & No & Temporary 3 & Caller-saved temporary \\
R16 & t4 & No & Temporary 4 & Caller-saved temporary \\
R17 & t5 & No & Temporary 5 & Caller-saved temporary \\
R18 & t6 & No & Temporary 6 & Caller-saved temporary \\
R19 & t7 & No & Temporary 7 & Caller-saved temporary \\
R20 & t8 & No & Temporary 8 & Caller-saved temporary \\
R21 & t9 & No & Temporary 9 & Caller-saved temporary \\
R22 & t10 & No & Temporary 10 & Caller-saved temporary \\
R23 & t11 & No & Temporary 11 & Caller-saved temporary \\
R24 & t12 & No & Temporary 12 & Caller-saved temporary \\
R25 & t13 & No & Temporary 13 & Caller-saved temporary \\
R26 & t14 & No & Temporary 14 & Caller-saved temporary \\
R27 & t15 & No & Temporary 15 & Caller-saved temporary \\
R28 & t16 & No & Temporary 16 & Caller-saved temporary \\
R29 & t17 & No & Temporary 17 & Caller-saved temporary \\
R30 & t18 & No & Temporary 18 & Caller-saved temporary \\
R31 & ra & No & Return address & Set by CALL, used by RET \\
\hline
\end{tabularx}
\end{table}

\section{Calling Convention}

\subsection{Function Prologue}

\begin{lstlisting}[style=assembly, caption={Standard Function Prologue}]
function_name:
    PUSH R31            ; Save return address
    PUSH R3             ; Save old stack pointer (optional)
    SUBI R3, R3, N      ; Allocate N bytes for locals
    ; Save callee-saved registers if used
\end{lstlisting}

\subsection{Function Epilogue}

\begin{lstlisting}[style=assembly, caption={Standard Function Epilogue}]
    ; Restore callee-saved registers
    ADDI R3, R3, N      ; Deallocate local variables
    POP R3              ; Restore stack pointer (if saved)
    POP R31             ; Restore return address
    RET                 ; Return to caller
\end{lstlisting}

\subsection{Argument Passing}

\begin{itemize}
    \item Arguments 1-8: Passed in R4-R11
    \item Additional arguments: Passed on stack
    \item Return value: R1
    \item Return address: R31 (set by CALL)
\end{itemize}

% ============================================================================
% APPENDIX C: SYSTEM CALL REFERENCE
% ============================================================================
\chapter{System Call Reference}

\section{System Call Interface}

System calls use the \code{SYSCALL} instruction with the following convention:

\begin{itemize}
    \item \textbf{R1}: System call number
    \item \textbf{R2}: First argument
    \item \textbf{R3}: Second argument  
    \item \textbf{R4}: Third argument
    \item \textbf{R1}: Return value (after syscall)
\end{itemize}

\section{System Call Table}

\begin{table}[H]
\centering
\caption{A1M System Call Numbers and Signatures}
\label{tab:syscall-reference}
\small
\begin{tabularx}{\textwidth}{|l|l|X|X|}
\hline
\textbf{Num} & \textbf{Name} & \textbf{Arguments} & \textbf{Returns} \\
\hline
0 & exit & R2: exit\_code & Does not return \\
\hline
1 & write & R2: fd, R3: buf, R4: len & R1: bytes written \\
\hline
2 & read & R2: fd, R3: buf, R4: len & R1: bytes read \\
\hline
3 & open & R2: path, R3: flags & R1: file descriptor \\
\hline
4 & close & R2: fd & R1: 0 on success \\
\hline
5 & malloc & R2: size & R1: address \\
\hline
6 & free & R2: address & R1: 0 on success \\
\hline
10 & print\_int & R2: value & - \\
\hline
11 & print\_str & R2: string\_addr & - \\
\hline
12 & read\_int & - & R1: integer value \\
\hline
13 & read\_str & R2: buffer, R3: max\_len & R1: bytes read \\
\hline
14 & print\_char & R2: character & - \\
\hline
15 & read\_char & - & R1: character \\
\hline
\end{tabularx}
\end{table}

\section{System Call Examples}

\subsection{Exit Program}

\begin{lstlisting}[style=assembly]
    MOVI R1, 0      ; syscall: exit
    MOVI R2, 0      ; exit code: 0 (success)
    SYSCALL 0
\end{lstlisting}

\subsection{Print Integer}

\begin{lstlisting}[style=assembly]
    MOVI R1, 10     ; syscall: print_int
    MOVI R2, 42     ; value to print
    SYSCALL 0
\end{lstlisting}

\subsection{Read and Write File}

\begin{lstlisting}[style=assembly]
.section .data
filename: .string "data.txt"
buffer:   .space 1024

.section .text
    ; Open file
    MOVI R1, 3          ; syscall: open
    MOVI R2, filename   ; file path
    MOVI R3, 0          ; flags: O_RDONLY
    SYSCALL 0
    MOV R10, R1         ; Save file descriptor
    
    ; Read from file
    MOVI R1, 2          ; syscall: read
    MOV R2, R10         ; file descriptor
    MOVI R3, buffer     ; buffer address
    MOVI R4, 1024       ; max bytes to read
    SYSCALL 0
    MOV R11, R1         ; Save bytes read
    
    ; Write to stdout
    MOVI R1, 1          ; syscall: write
    MOVI R2, 1          ; fd: stdout
    MOVI R3, buffer     ; buffer address
    MOV R4, R11         ; bytes to write
    SYSCALL 0
    
    ; Close file
    MOVI R1, 4          ; syscall: close
    MOV R2, R10         ; file descriptor
    SYSCALL 0
\end{lstlisting}

% ============================================================================
% APPENDIX D: ALGO LANGUAGE GRAMMAR
% ============================================================================
\chapter{Algo Language Complete Grammar}

\section{EBNF Notation}

\begin{itemize}
    \item \code{::=} - Production rule
    \item \code{|} - Alternative
    \item \code{*} - Zero or more
    \item \code{+} - One or more
    \item \code{?} - Optional (zero or one)
    \item \code{( )} - Grouping
    \item \code{UPPERCASE} - Terminal (keyword or token type)
    \item \code{lowercase} - Non-terminal
\end{itemize}

\section{Complete Grammar}

\begin{verbatim}
program         ::= declaration*

declaration     ::= function_decl
                  | var_decl

function_decl   ::= 'fun' IDENTIFIER '(' parameter_list? ')' 
                    ('->' type)? block

parameter_list  ::= parameter (',' parameter)*

parameter       ::= type IDENTIFIER ('[' ']')?

var_decl        ::= type IDENTIFIER ('[' INTEGER ']')? 
                    ('=' expression)? ';'

block           ::= '{' statement* '}'

statement       ::= var_decl
                  | expr_stmt
                  | if_stmt
                  | while_stmt
                  | for_stmt
                  | return_stmt
                  | block

expr_stmt       ::= expression ';'

if_stmt         ::= 'if' '(' expression ')' statement
                    ('else' statement)?

while_stmt      ::= 'while' '(' expression ')' statement

for_stmt        ::= 'for' '(' (var_decl | expr_stmt | ';')
                    expression? ';'
                    expression? ')' statement

return_stmt     ::= 'return' expression? ';'

expression      ::= assignment

assignment      ::= (IDENTIFIER ('[' expression ']')?)
                    '=' assignment
                  | logic_or

logic_or        ::= logic_and ('||' logic_and)*

logic_and       ::= equality ('&&' equality)*

equality        ::= comparison (('==' | '!=') comparison)*

comparison      ::= term (('<' | '>' | '<=' | '>=') term)*

term            ::= factor (('+' | '-') factor)*

factor          ::= unary (('*' | '/' | '%') unary)*

unary           ::= ('!' | '-') unary
                  | postfix

postfix         ::= primary ('[' expression ']'
                  | '(' argument_list? ')')*

primary         ::= INTEGER
                  | STRING
                  | CHAR
                  | 'true'
                  | 'false'
                  | IDENTIFIER
                  | '(' expression ')'

argument_list   ::= expression (',' expression)*

type            ::= 'int' | 'long' | 'char' | 'bool'
\end{verbatim}

\section{Operator Precedence}

\begin{table}[H]
\centering
\caption{Algo Operator Precedence (Highest to Lowest)}
\label{tab:operator-precedence}
\begin{tabularx}{\textwidth}{|l|X|l|}
\hline
\textbf{Precedence} & \textbf{Operators} & \textbf{Associativity} \\
\hline
1 (Highest) & Function call \code{()}, Array access \code{[]} & Left \\
\hline
2 & Unary \code{!}, \code{-} & Right \\
\hline
3 & Multiplicative \code{*}, \code{/}, \code{\%} & Left \\
\hline
4 & Additive \code{+}, \code{-} & Left \\
\hline
5 & Relational \code{<}, \code{>}, \code{<=}, \code{>=} & Left \\
\hline
6 & Equality \code{==}, \code{!=} & Left \\
\hline
7 & Logical AND \code{\&\&} & Left \\
\hline
8 & Logical OR \code{||} & Left \\
\hline
9 (Lowest) & Assignment \code{=} & Right \\
\hline
\end{tabularx}
\end{table}

\section{Type System}

\subsection{Primitive Types}

\begin{table}[H]
\centering
\caption{Algo Primitive Types}
\begin{tabularx}{\textwidth}{|l|l|l|X|}
\hline
\textbf{Type} & \textbf{Size} & \textbf{Range} & \textbf{Description} \\
\hline
int & 64 bits & $-2^{63}$ to $2^{63}-1$ & Signed integer \\
\hline
long & 64 bits & $-2^{63}$ to $2^{63}-1$ & Alias for int \\
\hline
char & 8 bits & 0 to 255 & Single character (ASCII) \\
\hline
bool & 8 bits & 0 or 1 & Boolean (true/false) \\
\hline
\end{tabularx}
\end{table}

\subsection{Array Types}

Arrays are declared with fixed sizes:

\begin{verbatim}
int arr[10];        // Array of 10 integers
char str[256];      // Character array (string)
bool flags[8];      // Array of booleans
\end{verbatim}

% ============================================================================
% APPENDIX E: ERROR MESSAGES
% ============================================================================
\chapter{Compiler Error Messages}

\section{Lexical Errors}

\begin{table}[H]
\centering
\caption{Lexer Error Messages}
\begin{tabularx}{\textwidth}{|l|X|}
\hline
\textbf{Error Code} & \textbf{Message} \\
\hline
LEX001 & Unterminated string literal \\
LEX002 & Invalid character in input \\
LEX003 & Integer literal too large \\
LEX004 & Invalid escape sequence in string \\
LEX005 & Unterminated comment \\
\hline
\end{tabularx}
\end{table}

\section{Syntax Errors}

\begin{table}[H]
\centering
\caption{Parser Error Messages}
\begin{tabularx}{\textwidth}{|l|X|}
\hline
\textbf{Error Code} & \textbf{Message} \\
\hline
SYN001 & Expected token X, found Y \\
SYN002 & Unexpected end of file \\
SYN003 & Missing semicolon \\
SYN004 & Unmatched parenthesis/bracket/brace \\
SYN005 & Invalid function declaration \\
SYN006 & Invalid expression \\
\hline
\end{tabularx}
\end{table}

\section{Semantic Errors}

\begin{table}[H]
\centering
\caption{Semantic Error Messages}
\begin{tabularx}{\textwidth}{|l|X|}
\hline
\textbf{Error Code} & \textbf{Message} \\
\hline
SEM001 & Undefined variable: X \\
SEM002 & Duplicate variable declaration: X \\
SEM003 & Type mismatch: expected X, got Y \\
SEM004 & Function X not defined \\
SEM005 & Argument count mismatch for function X \\
SEM006 & Cannot assign to non-lvalue \\
SEM007 & Array index must be integer \\
SEM008 & Return type mismatch \\
\hline
\end{tabularx}
\end{table}

% ============================================================================
% BIBLIOGRAPHY
% ============================================================================
\begin{thebibliography}{99}

\bibitem{hennessy2017}
J. L. Hennessy and D. A. Patterson,
\textit{Computer Architecture: A Quantitative Approach},
6th ed. Morgan Kaufmann, 2017.

\bibitem{patterson2013}
D. A. Patterson and J. L. Hennessy,
\textit{Computer Organization and Design: The Hardware/Software Interface},
5th ed. Morgan Kaufmann, 2013.

\bibitem{aho2006}
A. V. Aho, M. S. Lam, R. Sethi, and J. D. Ullman,
\textit{Compilers: Principles, Techniques, and Tools},
2nd ed. Addison-Wesley, 2006.

\bibitem{muchnick1997}
S. S. Muchnick,
\textit{Advanced Compiler Design and Implementation},
Morgan Kaufmann, 1997.

\bibitem{waterman2016}
A. Waterman and K. Asanović, Eds.,
\textit{The RISC-V Instruction Set Manual, Volume I: User-Level ISA},
RISC-V Foundation, 2016.

\bibitem{sweetman2006}
D. Sweetman,
\textit{See MIPS Run},
2nd ed. Morgan Kaufmann, 2006.

\bibitem{thomas2009}
D. E. Thomas and P. R. Moorby,
\textit{The Verilog Hardware Description Language},
5th ed. Springer, 2009.

\bibitem{harris2012}
S. L. Harris and D. M. Harris,
\textit{Digital Design and Computer Architecture},
2nd ed. Morgan Kaufmann, 2012.

\end{thebibliography}

\end{document}
